
Let $\C$ be an $\infty$-category. Recall that a full subcategory $\iota\colon \D \subseteq \C$ is called reflective if the inclusion functor $\iota$ admits a left adjoint $L\colon \C \to \D$. In this case we say that $L$ exhibits $\D$ as a left Bousfield localisation of $\C$. In particular, $L$ is a localisation functor: it exhibits $\D$ as the localisation of $\C$ by the collection of $L$-equivalences, that is, the maps $f\colon x \to y$ in $\C$ such that $L(f)$ is an equivalence in $\D$.
In this situation, there is well-known machinery to transport a symmetric monoidal structure on $\C$, or more generally an $\cO$-monoidal structure on $\C$ over some $\infty$-operad $\cO$, to $\D$, 
see e.g.\ \cite{HA}*{\S 2.2.1}. The goal of the present appendix is to treat the more general case where $\C$ is only required to be an $\infty$-operad (over $\cO$), and discuss left Bousfield localisations in this setting, as we introduce in \S\reffour{subsection:multiplicative-localisation} under the name of $\cO$-multiplicative localisations. A special type of multiplicative localisations is given by the smashing localisations associated to idempotent objects, which can be defined when $\C$ and $\cO$ admit compatible unit objects. We develop the relevant notions and study smashing localisations in this generality in \S\reffour{subsection:units-sym-mon-objects} and 
\S\reffour{subsection:smashing-localisations}.


The application we have in mind is the following. Suppose given a small symmetric monoidal $\infty$-category and a presentably symmetric monoidal $\infty$-category $\E$. Then the functor $\infty$-category $\Fun(\C,\E)$ inherits a symmetric monoidal structure, known as Day convolution. Given a collection $\cR$ of cone diagrams in $\C$, one can consider the full subcategory $\Fun^\cR(\C,\E) \subseteq \Fun(\C,\E)$ spanned by those functors taking all diagrams in $\cR$ to limit diagrams in $\E$. The adjoint functor theorem then implies that this inclusion has a left adjoint, and there are well-known criteria for this adjoint to be compatible with Day convolution (see, e.g., Lemma~\refone{lemma:day-localisation}). %
If we now relax the condition that $\C$ is small and only assume that $\C$ is, say, locally small, then one can still define the Day convolution $\infty$-operad $\Fun(\C,\E)^\Day$, which is generally not a symmetric monoidal $\infty$-category. In addition, given a collection of cone diagrams $\cR$, a left adjoint to the inclusion $\Fun^\cR(\C,\E) \subseteq \Fun(\C,\E)$ might only be defined on some intermediate subcategory. %
As a test case of interest, we have dedicated \S\reffour{subsection:day-convolution} to a discussion of multiplicative and smashing localisations in the setting of Day convolution $\infty$-operads. The results we obtain here play a crucial role in the body of the paper for establishing universal lax symmetric monoidal structures on Grothendieck-Witt and L-theory; we refer the reader to the beginning of \S\reffour{subsection:multiplicative-GW} for a discussion of exactly which results are used and how.


\subsection{Multiplicative Bousfield localisations}
\label{subsection:multiplicative-localisation}%



We begin with some basic definitions and notations. Following Lurie's conventions, for an $\infty$-operad $\cO^\otimes$, we write $\cO := \cO^\otimes_{\langle 1 \rangle}$ for its $\infty$-category of colours and refer to maps in $\cO^{\otimes}$ with target in $\cO$ as \defi{operations}. Given $p\colon\C^\otimes \rightarrow \cO^\otimes$ a map of $\infty$-operads, we say that $p$ exhibits $\C^\otimes$ as \defi{$\cO$-monoidal} if $p$ is a cocartesian fibration of $\infty$-operads. A functor $f^\otimes \colon \C^\otimes \to \D^\otimes$ over $\cO^\otimes$ between $\cO$-monoidal $\infty$-categories is called $\cO$-monoidal if it is a map of $\infty$-operads which preserves cocartesian arrows over $\cO^{\otimes}$. We also recall that the maps in any $\infty$-operad support a factorisation system into \defi{active} and \defi{inert} maps, see \cite{HA}*{\S 2.1} for the definition and a thorough discussion. 



\begin{definition}
\label{definition:left-bousfield}%
We say that a map $L^{\otimes}$ as in the following diagram
\[
\begin{tikzcd}
\C^{\otimes}\ar[rr, "L^{\otimes}"']\ar[dr,"p"'] && \D^{\otimes}\ar[dl,"q"]\ar[ll,dotted, bend right = 20, "R^{\otimes}"',"\rotatebox{180}{$\perp$}"] \\
& \cO^{\otimes} &
\end{tikzcd}
\]
is 
an \defi{$\cO$-multiplicative} left Bousfield localisation if it admits a
right adjoint $R^{\otimes}$ relative to $\cO^{\otimes}$ (see~\cite{HA}*{Definition 7.3.2.2}) which is furthermore %
fully faithful and preserves inert maps. %

If in addition $L^{\otimes}$ is an $\cO$-monoidal functor among $\cO$-monoidal $\infty$-categories then we say that $L^{\otimes}$ is an \defi{$\cO$-monoidal} left Bousfield localisation.
\end{definition}



\begin{remark}
\label{remark:automatically-monoidal}%
If $L^{\otimes}\colon \C^{\otimes} \to \D^{\otimes}$ is an $\cO$-multiplicative left Bousfield localisation and $\C^{\otimes}$ is $\cO$-monoidal then $\D^{\otimes}$ is also $\cO$-monoidal by \cite{NS}*{Proposition A.14} and $L^{\otimes}$ preserves cocartesian arrows over $\cO^{\otimes}$ since it is a relative left adjoint and hence preserves relative colimits. In particular, an $\cO$-multiplicative left Bousfield localisation is $\cO$-monoidal as soon as its domain is $\cO$-monoidal.
\end{remark}




\begin{remark}
\label{remark:fibrewise-L}%
In the situation of Definition~\reffour{definition:left-bousfield},
the left adjoint $L^{\otimes}$ is also a fibrewise left adjoint to $R^{\otimes}\colon \D^{\otimes} \hrar \C^{\otimes}$ by~\cite{HA}*{Proposition 7.3.2.5}, and hence determines an adjunction $L_a\colon \C_a \adj \D_a \cocolon R_a$ for every colour $a \in \cO$. In addition, for every $A=(a_1,\dots,a_n) \in \cO^{\otimes}$ 
the induced adjunction on fibres over $A$ breaks as a pointwise adjunction of cartesian products
\[
\prod L_{a_i} \colon \prod \C_{a_i} \adj \prod \D_{a_i} \cocolon \prod R_{a_i},
\]
since $L^{\otimes}$ and $R^{\otimes}$ are maps of $\infty$-operads. 
\end{remark}

In the situation of Definition~\reffour{definition:left-bousfield},
the counit of the adjunction $\C^{\otimes} \adj \D^{\otimes}$ is an equivalence. It hence induces, for every $\infty$-operad $\E^{\otimes} \to \cO^{\otimes}$ over $\cO^{\otimes}$, an adjunction
\begin{equation}
\label{eq:adj-algberas}%
\Alg_{\C/\cO}(\E) \locadj \Alg_{\D/\cO}(\E) 
\end{equation}
on $\infty$-categories of algebras, whose left and right adjoints are given by restriction along $R^{\otimes}\colon \D^{\otimes} \hookrightarrow \C^{\otimes}$ and $L^{\otimes}\colon \C^{\otimes} \to \D^{\otimes}$, respectively, and whose counit is again an equivalence. The essential image of the resulting fully faithful functor 
\[
(L^{\otimes})^*\colon \Alg_{\D/\cO}(\E) \to \Alg_{\C/\cO}(\E)
\]
then consists of those $\infty$-operad maps $\Psi^{\otimes}\colon \C^{\otimes} \to \E^{\otimes}$ over $\cO$ which send the components of the unit of $\C^{\otimes} \locadj \D^{\otimes}$ to equivalences. By Remark~\reffour{remark:fibrewise-L}, this is the same as saying that $\Psi_a\colon \C_a \to \E_a$ sends the components of the unit of $\C_a \locadj \D_a$ to equivalences for every $a \in \cO$, or equivalently, that $\Psi_a$ inverts $L_a$-equivalences for every $a \in \cO$. We conclude:


\begin{corollary}
Let $L^{\otimes} \colon \C^{\otimes} \to \D^{\otimes}$ be an $\cO$-multiplicative left Bousfield localisation.
Then $L^{\otimes}$ exhibits $\D^{\otimes}$ as the initial recipient of an $\infty$-operad map over $\cO^{\otimes}$ from $\C^{\otimes}$ which inverts $L_a$-equivalences for every $a \in \cO$. %
\end{corollary}


\begin{lemma}
\label{lem:monoidal-bousfield}%
Let $L^{\otimes}\colon \C^{\otimes} \to \D^{\otimes}$ be an $\cO$-monoidal left Bousfield localisation and $\E^{\otimes} \to \cO^{\otimes}$ an $\cO$-monoidal $\infty$-category.
Then the adjunction~\eqreffour{eq:adj-algberas}
restricts to an equivalence 
\[
\Fun^{\cO-\otimes}_L(\C^{\otimes},\E^{\otimes}) \locadjL{\simeq} \Fun^{\cO-\otimes}(\D^{\otimes},\E^{\otimes})
\]
where on the left hand side we have the full subcategory 
spanned by those $\cO$-monoidal functors $\Psi^{\otimes}\colon \C^{\otimes} \to \E^{\otimes}$ 
which invert $L_a$-equivalences 
for every $a \in \cO$.
\end{lemma}
\begin{proof}
Note first that since $L^{\otimes}$ is an $\cO$-monoidal functor which inverts $L_a$-equivalences for every $a \in \cO$ the same holds for $L^{\otimes} \circ \Phi^{\otimes}$ for every $\cO$-monoidal functor $\Phi^\otimes \colon \D^{\otimes} \to \E^{\otimes}$.
Moreover, if $e\colon (x_1,...,x_n) \to y$ is a cocartesian arrow in $\D^{\otimes}$ over an arrow of the form $(a_1,...,a_n) \to b$ in $\cO^{\otimes}$, then its image in $\C^{\otimes}$ factors as a composite of a cocartesian arrow $(x_1,...,x_n) \to y'$ in $\C^{\otimes}$ over $\cO^{\otimes}$ followed by an $L_b$-equivalence $y'\to y$. It then follows that if $\Psi^{\otimes} \colon  \C^{\otimes} \to \E^{\otimes}$ inverts $L_a$-equivalences for every $a \in \cO$ then the composite $\Psi^{\otimes} \circ \iota^{\otimes}$ is an $\cO$-monoidal functor. This shows that the adjunction~\eqreffour{eq:adj-algberas}
indeed restricts as stated. At the same time, the counit of this adjunction is an equivalence already before restriction, and the unit is an equivalence after restriction since its components are equivalences when evaluated on any $\infty$-operad map that inverts $L_a$-equivalences. 
\end{proof}





We now come to the main result of this subsection, whose goal is to establish a recognition criterion for $\cO$-multiplicative left Bousfield localisations. Our result is a direct generalisation of \cite{HA}*{Proposition 2.2.1.9}, where the analogous criterion is established in the $\cO$-monoidal case.


\begin{definition}
\label{definition:multiplicative-localisation}%
Let $\C^{\otimes} \to \cO^{\otimes}$ be a map of $\infty$-operads. %
We will say that a collection of left Bousfield localisations $\{L_a\colon \C_a \to \D_a \subseteq \C_a\}_{a \in \cO}$ is \defi{$\cO$-multiplicative} if for every operation $\alp\colon (a_1,...,a_n) \to b$ in $\cO$, any choice of an $L_{a_i}$-equivalence $f_i\colon x_i \to y_i$ in $\C_{a_i}$ for each $i=1,...,n$, and any $d \in \D_b$, restriction along the $f_i$ induces an equivalence
\[
\Mul^{\alp}_{\C}(y_1,\dots,y_n;d) \xrightarrow{\simeq} \Mul^{\alp}_{\C}(x_1,\dots,x_n; d)
\]
on spaces of operations in $\C$ covering $\alp$. If $\cO = \Com$ then we will drop $\cO$ from the notation and just write \defi{multiplicative}.
\end{definition}

\begin{remarks}\
\label{remark:full-suboperad}%
\begin{enumerate}
\item If $\C^{\otimes} \to \cO^{\otimes}$ is an $\cO$-monoidal $\infty$-category, 
then the condition that a given collection of localisations $\{L_a\}_{a \in \cO}$ is $\cO$-multiplicative can be rephrased by saying that for every operation $\alp\colon (a_1,...,a_n) \to b$ and every $i = 1,...,n$, the induced map $\C_{a_1} \times \dots \times \C_{a_n} \to \C_{b}$ sends $L_{a_i}$-equivalences in the $i$th coordinate to $L_{b}$-equivalences. In particular, we recover in this case \cite{HA}*{Definition 2.2.1.6}.
\item
\label{item:full-suboperad}%
It follows from the definition that if $\C^{\otimes}$ is an $\infty$-operad over $\cO^{\otimes}$ equipped with an $\cO$-multiplicative collection of left Bousfield localisations $\{L_a\colon \C_a \to \D_a \subseteq \C_a\}_{a \in \cO}$ and $\E^{\otimes} \subseteq \C^{\otimes}$ is a full suboperad such that $\E_a$ contains $\D_a$ for every $a \in \cO$ then the collection of restricted localisations $\{L_a|_{\E_a}\colon \E_a \to \D_a \subseteq \E_a\}_{a \in \cO}$ is again $\cO$-multiplicative.
\item
\label{item:one-at-a-time}%
In the situation of Definition~\reffour{definition:multiplicative-localisation}, one can break the induced map on spaces of operations as a composite where each step features just one of the $f_i$, with the rest of the maps being identities. In establishing multiplicativity one may therefore assume that all the $f_i$'s are identities except for a single $f_{i_0}$. In addition, up to replacing $\alp$ with one of its permutations we may as assume that $i_0=1$.
\end{enumerate}
\end{remarks}






\begin{proposition}
\label{proposition:multiplicative-localisation}%
Let $p\colon\C^{\otimes} \to \cO^{\otimes}$ be map of $\infty$-operads and $\{L_a\colon \C_a \to \D_a\}_{a \in \cO}$ an $\cO$-multiplicative collection of left Bousfield localisations as in Definition~\reffour{definition:multiplicative-localisation}. 
Then the inclusion $\D^{\otimes} \subseteq \C^{\otimes}$ of the full suboperad spanned by the $\D_a$ admits a left adjoint
\[
\begin{tikzcd}
\C^{\otimes} \ar[rr,"L^{\otimes}"]\ar[dr,"p"'] && \D^{\otimes} \ar[dl,"q"] \\
& \cO^{\otimes} & 
\end{tikzcd}
\]
relative to $\cO^{\otimes}$, and such that $L^{\otimes}$ preserves inert arrows. %
In particular, $L^{\otimes}$ is an $\cO$-multiplicative left Bousfield localisation.
\end{proposition}
\begin{proof}
We show that the inclusion $\D^{\otimes} \subseteq \C^{\otimes}$ admits a left adjoint $L^\otimes$ relative to $\cO^{\otimes}$. Then $L^\otimes$ preserves colimits relative to $\cO^\otimes$, and hence in particular $\cO^\otimes$-cocartesian arrows. In particular, $L^\otimes$ is a map of $\infty$-operads.

Applying~\cite{HA}*{Proposition 7.3.2.1} it will suffice to show that for every $\uline{x} := (x_1,\dots,x_n) \in \C^{\otimes}$, the slice category $\D^{\otimes} \times_{\C^{\otimes}} \C^{\otimes}_{\uline{x}/}$ admits an initial object $f\colon \uline{x}\to \uline{d}$ satisfying the additional property that the image of $f$ in $\cO^{\otimes}$ is an equivalence. Let $p(\uline{x}) := \uline{a} = (a_1,\dots,a_n)$. Since each inclusion $\D_{a_i} \subseteq \C_{a_i}$ admits a left adjoint $L_{a_i}\colon \C_{a_i} \to \D_{a_i}$, each of the slice categories $\D_{a_i} \times_{\C_{a_i}} (\C_{a_i})_{x_i/}$ admits an initial object $f_i\colon x_i \to d_i$. Let $f\colon \uline{x}=(x_i,\dots,x_n) \to \uline{d}=(d_1,\dots,d_n)$ be the induced map, so that $\uline{d} \in \D^{\otimes}$ and $f$ maps to the identity on $\uline{a}$ by construction. We now claim that $f$ encodes an initial object of $\D^{\otimes} \times_{\C^{\otimes}} \C^{\otimes}_{\uline{x}/}$. 
For this, we need to show that for every $\uline{e}=(e_1,\dots,e_m) \in \D^{\otimes}$ the pre-composition map
\[
(-) \circ f\Map_{\C^{\otimes}}(\uline{d},\uline{e}) \to \Map_{\C^{\otimes}}(\uline{x},\uline{e})
\]
is an equivalence. Writing $b_i=f(e_i)$ and $\uline{b}=(b_1,\dots,b_m)$, the above map fits in a commutative diagram
\[
\begin{tikzcd}
\Map_{\C^{\otimes}}(\uline{d},\uline{e}) \ar[rr]\ar[dr] && \Map_{\C^{\otimes}}(\uline{x},\uline{e})\ar[dl] \\
&\Map_{\cO^{\otimes}}(\uline{a},\uline{b}) &
\end{tikzcd}
\]
It will then suffice to show that the top horizontal map induces an equivalence on vertical fibres for every map $\alp\colon \uline{a} \to \uline{b}$ in $\cO^{\otimes}$. Let $\rho\colon \langle n\rangle \to \langle m\rangle$ be the image of $\alp$ in $\Com$. Then $\alp$ determines a map $\alp_j\colon (a_i \;|\; \rho(i)=j) \to b_j$ for each $j=1,...,m$ and the mapping spaces in the top row decompose compatibly as products of mapping spaces covering the various $\alp_i$. We may hence assume without loss of generality that $m=1$. Then, we need to show that for any operation 
$\alp\colon \uline{a} \to b$ and any $e \in \D_b$, the induced map on spaces of operations in $\C^\otimes$ covering $\alp$
\[
\Mul^{\alp}_{\C}(d_1,\dots,d_n;e) \to \Mul^{\alp}_{\C}(x_1,\dots,x_n;e)
\]
is an equivalence, which is the case by the assumption that the collection of localisations $\{L_a\}_{a \in \cO}$ is $\cO$-multiplicative.
\end{proof}

When $\C^{\otimes} \to \cO^{\otimes}$ is an $\cO$-monoidal $\infty$-category Proposition~\reffour{proposition:multiplicative-localisation} recovers \cite{HA}*{Proposition 2.2.1.9}, see Remark \reffour{remark:automatically-monoidal}, as we record now (we also record the corresponding universal property of $\D^{\otimes}$, which is not mentioned in loc.\ cit.)

\begin{proposition}
\label{proposition:multiplicative-localisation-2}%
In the situation of Proposition~\reffour{proposition:multiplicative-localisation}, suppose in addition that $\C^{\otimes} \to \cO^{\otimes}$ is an $\cO$-monoidal $\infty$-category. Then $\D^{\otimes} \to \cO^{\otimes}$ is also $\cO$-monoidal as is the functor $L^{\otimes}$ provided by Proposition~\reffour{proposition:multiplicative-localisation}. In addition, this functor
enjoys the following universal property: for every $\cO$-monoidal $\infty$-category $\E^{\otimes} \to \cO^{\otimes}$, the functor
\[
(L^{\otimes})^*\colon \Fun^{\cO-\otimes}(\D^{\otimes},\E^{\otimes}) \to \Fun^{\cO-\otimes}(\C^{\otimes},\E^{\otimes})
\]
is fully faithful, and its essential image is spanned by those $\cO$-monoidal functors $\Psi^{\otimes}\colon \C^{\otimes} \to \E^{\otimes}$ such that $\Psi_a\colon \C_a \to \D_a$ sends $L_a$-equivalences to equivalences for every $a \in \cO$.
\end{proposition}
\begin{proof}
Combine Remark \reffour{remark:automatically-monoidal} and Lemma~\reffour{lem:monoidal-bousfield}.
\end{proof}


\begin{example}
\label{example:sm-localisation}%
Suppose $\cO = \Com$. Then an $\cO$-monoidal $\infty$-category amounts to a symmetric monoidal $\infty$-category $\C^{\otimes}$, and a left localisation $L\colon \C \to \D \subseteq \C$ is multiplicative if and only if for every pair of $L$-equivalences $f\colon x \to y$ and $g\colon z \to w$, the map $f \otimes g \colon x \otimes z \to y \otimes w$ is again an $L$-equivalence, equivalently, the collection of $L$-equivalences is closed under the operation $x \otimes (-)\colon \C \to \C$ for every $x \in \C$. 
When this holds, the symmetric monoidal structure of $\C$ descends to one on $\D$, and the resulting symmetric monoidal functor
\[
L^{\otimes} \colon \C^{\otimes} \to \D^{\otimes}
\]
is initial in symmetric monoidal functors out of $\C^{\otimes}$  which send $L$-equivalences to equivalence, as an instance of Proposition~\reffour{proposition:multiplicative-localisation}.
The induced symmetric monoidal structure on $\D$ is then given explicitly by the formula $x \otimes y = L(\iota(x) \otimes \iota(y))$.
\end{example}

\begin{example}
\label{example:mcom-loc}%
Suppose $\cO=\MCom$, the $\infty$-operad which classifies pairs $(R,M)$ of a commutative algebra object $R$ equipped with an $R$-module $M$, see~\cite{hinich-modules}*{\S 5.2}. Then $\MCom$-monoidal $\infty$-categories correspond to pairs $(\C^{\otimes},\M)$ where $\C$ is a symmetric monoidal $\infty$-category and $\M$ is an $\infty$-category equipped with an action of $\C$. A pair of left localisations $L_{\alg}\colon \C \to \D \subseteq \C$ and $L_{\md}\colon \M \to \cN \subseteq \M$ is then $\MCom$-multiplicative if $L_{\alg}$ is a multiplicative localisation of $\C$ as in Example~\reffour{example:sm-localisation} and the action of $\C$ on $\M$ satisfies that for $f\colon x \to y$ in $\C$ and $g\colon z \to w$ in $\M$, if $f$ is an $L_{\alg}$-equivalence and $g$ is an $L_{\md}$-equivalence then $f \otimes g \colon x \otimes z \to y \otimes w$ is an $L_{\md}$-equivalences. In that case we will also say that the localisation $L_{\md}$ is \emph{linear over $L_{\alg}$}. By Proposition~\reffour{proposition:multiplicative-localisation}, when this holds the symmetric monoidal structure of $\C$ descends to one on $\D$ and the $\C$-action on $\M$ descends to a compatible $\D$-action on $\cN$, such that the resulting map
\[
L^{\otimes}\colon (\C^{\otimes},\M) \to (\D^{\otimes},\cN)
\]
of $\MCom$-algebras is initial in $\MCom$-algebras maps out of $(\C,\M)$ sending $L_{\alg}$-equivalences and $L_{\md}$-equivalences to equivalences. The induced $\D$ action on $\cN$ is given explicitly by the formula $x \otimes y = L_{\md}(\iota_{\alg}(x) \otimes \iota_{\md}(y))$, where $\iota_{\alg}\colon \D \subseteq \C$ and $\iota_{\md}\colon \cN \subseteq \M$ are the corresponding full inclusions.

An interesting special case of the above is when $\D=\C$ and $L_{\alg} = \id$.
In this case we can ignore the algebra component and just consider the left localisation $L\colon \M \to \cN \subseteq \M$ on $\M$. To check multiplicativity it is then enough to check that the collection of $L$-equivalences in $\M$ is preserved by the operation $x \otimes (-)\colon \M \to \M$ for every $x \in \C$. In this case we will also say that $L$ is a \defi{$\C$-linear localisation}. Then $\cN$ inherits from $\M$ its $\C$-module structure and the resulting $\C$-module refinement $\M \to \cN$ of $L$ is initial in $\C$-module functors out of $\M$ sending $L$-equivalences to equivalences.
\end{example}




\subsection{Operadic units and idempotents}
\label{subsection:units-sym-mon-objects}%




Our goal in this subsection is to recall and explore the notions of units and idempotent objects in $\infty$-operads which will be used in \S\reffour{subsection:smashing-localisations} for the study of smashing localisations of $\infty$-operads (with units). The notion of units we use here is closely related to that introduced by Lurie in~\cite{HA}*{\S 3.2.1}, see Remark~\reffour{remarks:symmetric-monoidal-objects-and-units}\;\reffouritem{item:HA}
for the precise relation between the two.

We begin with the auxiliary notion of an operadically cocartesian arrow. To that end, recall that for an $\infty$-operad $\cO^\otimes$ there is a symmetric monoidal structure on $\cO^\otimes_\act$ which we refer to as concatenation, see \cite{HA}*{Corollary 2.2.4.5 \& Remark 2.2.4.6}.  
Recall from~\cite{HA}*{Definition 3.1.1.2} that for a map of $\infty$-operads $p\colon \C^{\otimes} \to \cO^{\otimes}$, a cone diagram $\ovl{\alp} \colon K^{\triangleright} \to \C^{\otimes}_{\act}$ is said to be a weak operadic $p$-colimit cone if the map 
\[
\C^{\act}_{\ovl{\alp}/} \to \C^{\act}_{\alp/} \times_{\cO^{\act}_{p\alp/}} \cO^{\act}_{p\ovl{\alp}/}
\]
is a trivial Kan fibration, where $\alp = \ovl{\alp}|_{K}$ and the superscript $\act$ means that we restrict attention to cones involving only active maps and having tips lying over $\<1\> \in \Com^{\otimes}$. We then say that $\ovl{\alp}$ is an operadic $p$-colimit cone if it is a weak operadic $p$-colimit cone and remains so after concatenation with arbitrary objects in $\C^{\otimes}$.
In what follows we will implicitly identify $\Delta^1$ with $(\Delta^0)^\triangleright$, so that morphisms can be viewed as cone diagrams.

\begin{lemma}
\label{lemma:operadically}%
Let $p\colon \C^{\otimes} \to \cO^{\otimes}$ be a map of $\infty$-operads and $\alp\colon x \to y$ a map in $\C^{\otimes}_\act$ with $y \in \C := \C^{\otimes}_{\<1\>}$. Then the following conditions are equivalent:
\begin{enumerate}
\item
\label{item:def1}%
$\alp$ encodes an operadic $p$-colimit cone in $\C^{\otimes}$.
\item
\label{item:def2}%
$\alp$ is $p$-cocartesian and remains $p$-cocartesian after concatenation with arbitrary $z \in \C^{\otimes}$.
\end{enumerate}
\end{lemma}

\begin{definition}
\label{definition:operadically-cocartesian}%
Let $p\colon \C^{\otimes} \to \cO^{\otimes}$ be a map of $\infty$-operads and $\alp$ a map in $\C^{\otimes}_\act$ with target in $\C$. We will say that $\alp$ is \defi{operadically $p$-cocartesian} if it satisfies the equivalent conditions of Lemma~\reffour{lemma:operadically}.
\end{definition}

\begin{proof}[Proof of Lemma~\reffour{lemma:operadically}]
To see $\reffouritem{item:def2} \Rightarrow \reffouritem{item:def1}$, it suffices by definition to show that every $p$-cocartesian active arrow in $\C^{\otimes}$
is a weak operadic $p$-colimit diagram.
But such an arrow is also $p_{\act}$-cocartesian, where $p_{\act}\colon \C_{\act}^\otimes \to \cO_{\act}^\otimes$, and so this follows directly from the definition of weak operadic $p$-colimits.
To get $\reffouritem{item:def1} \Rightarrow \reffouritem{item:def2}$ we first note that what we have just proven implies that active equivalences, in particular identities, are operadic $p$-colimit cones. Then the claim follows rom \cite{HA}*{Prop.\ 3.1.1.10} using that an arrow is $p$-cocartesian if and only if it encodes a $p$-colimit diagram.
\end{proof}


\begin{definition}
\label{definition:monoidal-and-unit}%
Let 
\[
\begin{tikzcd}
\C^{\otimes} \ar[dr,"p"']\ar[rr,"q"] && \cO^{\otimes} \ar[dl,"r"] \\
& \Com^{\otimes} &
\end{tikzcd}
\]
be a diagram of $\infty$-operads.
\begin{enumerate}
\item 
A map $u\colon \emptyset \to \one_\cO$ in $\cO^\otimes$ lying over the unique map $\<0\> \to \<1\>$ in $\Com^{\otimes}$ is said to exhibit $\one_\cO$ as a \defi{unit object} of $\cO$ if it is operadically $r$-cocartesian. We say that a map $\tilde{u}\colon \emptyset \to \one_\C$ in $\C^{\otimes}$ exhibits $\one_\C$ as a \defi{unit over $\one_\cO$} if it exhibits $\one_\C$ as a unit in $\C$ and $u=q(\tilde{u})$.
\item
An object $x \in \cO$ is said to be \defi{symmetric monoidal} if for every $y \in \cO$ there exists an operadically $r$-cocartesian map $(x,y) \to x\otimes y$ (using suggestive notation) lifting the unique active map $\mu_2\colon \<2\> \to \<1\>$.
Given a symmetric monoidal object $x \in \cO$, we will say that an object $\tilde{x}$ in $\C$ lifting $x$ is \defi{symmetric monoidal over $x$} if $\tilde{x}$ is symmetric monoidal %
and for every operadically $p$-cocartesian map %
$\alp\colon (\tilde{x},\tilde{y}) \to \tilde{x}\otimes\tilde{y}$ %
covering $\mu_2$, the map $q(\alp)$ is operadically $r$-cocartesian. 
\end{enumerate}
\end{definition}


\begin{examples}\
\begin{enumerate}
\item
In the $\infty$-operad $\Com$, the unique map $\<0\> \to \<1\>$ exhibits the $\<1\>$ as a unit object. 
\item
If $\D^{\otimes}$ is a symmetric monoidal $\infty$-category then the unit $\one_{\D}$ of $\D^{\otimes}$ is also canonically a unit object in the sense of Definition~\reffour{definition:monoidal-and-unit}.
\item
In the $\infty$-operad $\MCom$ the unique null operation $\emptyset \to a$ exhibits the algebra colour $a$ as a unit object. 
\end{enumerate}
\end{examples}



\begin{lemma}
\label{lemma:absolute-is-relative}%
Let $q\colon \C^{\otimes} \to \cO^{\otimes}$ be a map of $\infty$-operads such that $\cO^{\otimes}$ admits a unit object $u\colon \emptyset \to \one_{\cO}$ and $\C^{\otimes}$ admits a unit $\wtl{u}\colon \emptyset \to \one_{\cO}$
over $\one_{\cO}$. Then the following holds:
\begin{enumerate}
\item
The map $\wtl{u}$
is an operadically $q$-cocartesian lift of $u$.
\item
Let $a \in \cO^{\otimes}$ be a symmetric monoidal object and $\alp\colon (a,b) \to a\otimes b$ an operadically cocartesian lift of the map $\<2\> \to \<1\>$ in $\Com^{\otimes}$. Then for any symmetric monoidal object in $x \in \C$ over $a$ and any lift $y \in \C$ of $b$ there exists an operadically $q$-cocartesian arrow $(x,y) \to x\otimes y$ lifting $\alp$.
\end{enumerate}
\end{lemma}
\begin{proof}
To prove the first claim, note that by definition both $\wtl{u}$ and $u=q(\wtl{u})$ are operadically cocartesian over $\Com^{\otimes}$, that is, they are cocartesian over $\Com^{\otimes}$ and remain so after concatenation with arbitrary objects. 
We conclude from~\cite{HTT}*{Proposition 2.4.1.3} that $\wtl{u}$ is $q$-cocartesian and remains $q$-cocartesian after concatenation with arbitrary objects. In other words, $\wtl{u}$ is operadically $q$-cocartesian.

For the second claim, let us fix $\alp\colon (a,b) \to a\otimes b$, an operadically cocartesian lift of the map $\<2\> \to \<1\>$ such that $a$ is symmetric monoidal and $x,y \in \C$ lifts of $a,b$ such that $x$ is symmetric monoidal over $a$. Then by definition there exists an operadically cocartesian arrow $\beta\colon (x,y) \to x\otimes y$ lift of $\<2\> \to \<1\>$ in $\C$ such that $q(\beta)$ is operadically cocatesian in $\cO$ over $\Com^{\otimes}$. We then have that $\alp$ and $q(\beta)$ are two cocartesian lifts of $\<2\> \to \<1\>$ with the same domain, and so there exists an equivalence $\tau\colon \alp\simeq q(\beta)$ as objects of $\cO^{\otimes}_{(a,b)/} \times_{\Com^{\otimes}_{\<2\>/}} \{\<2\>\to \<1\>\}$. Using $\tau$ we may hence view $\beta$ as a lift of $\alp$. Arguing as above we again conclude from~\cite{HTT}*{Proposition 2.4.1.3} that $\beta$ is operadically $q$-cocartesian, as claimed.
\end{proof}




\begin{remarks}\
\label{remarks:symmetric-monoidal-objects-and-units}%
\begin{enumerate}
\item
If $\cO^{\otimes}$ is any $\infty$-operad and $u\colon \emptyset \to \one_{\cO}$  is a map in $\cO^{\otimes}$ lying above $\<0\> \to \<1\>$, then $u$ exhibits $\one_{\cO}$ as a unit object if and only if it exhibits it as a unit object over the unit of $\Com$. %
\item
\label{item:HA}%
If follows from Lemma~\reffour{lemma:absolute-is-relative} that a map $\tilde{u}\colon \emptyset \to \one_\C$ in $\C^{\otimes}$ exhibits $\one_\C$ as a unit over $\one_\cO$ if and only if it exhibits $\one_{\C}$ as an $\one_{\cO}$-unit in the sense of~\cite{HA}*{Definition 3.2.1.1}. To avoid confusion, let us note that loc.\ cit. also considers the unit property for objects of $\C$ lying over arbitrary objects in $\cO$, a notion which has no counterpart in Definition~\reffour{definition:monoidal-and-unit}.
\item
\label{item:unit-is-monoidal}%
Any unit object is symmetric monoidal, and if $\C^{\otimes} \to \cO^{\otimes}$ is a map of $\infty$-operad such that $\cO^{\otimes}$ has a unit $\one_{\cO}$ then any unit in $\C$ over $\one_{\cO}$ is also symmetric monoidal over $\one_{\cO}$. %
\item
\label{item:simplifies}%
In the $\infty$-operad $\Com^{\otimes}$ the unique active arrow $\mu_n\colon \<n\> \to \<1\>$ with $n \geq 2$ can be written as a composite of arrows obtained from the active arrow $\mu_2\colon \<2\> \to \<1\>$ via concatenation with objects. Given an $\infty$-operad $\C^{\otimes}$ it then follows by induction that if $x \in \C$ is a symmetric monoidal object then for every $y \in \C$ there exists an operadically cocartesian map of the form $\alp\colon (x,\dots,x,y) \to x\otimes ... \otimes x\otimes y$ lying over $\mu_n$. Similarly, if $x \in \C$ is symmetric monoidal over some symmetric monoidal object in $\cO$ then the operadically cocartesian maps $\alp$ just mentioned are sent to operadically cocartesian arrows in $\cO^{\otimes}$ (over $\Com^{\otimes}$). %
\item
\label{item:local-unit}%
Suppose that $\C^{\otimes}$ admits a unit $\emptyset \to \one_{\C}$ over a unit $\one_{\cO}$ in $\cO$ and let $L^{\otimes} \colon \C^{\otimes} \to \D^{\otimes}$ be an $\cO$-multiplicative left Bousfield localisation. %
Then the composite 
\[
\emptyset \xrightarrow{u} \one_{\C} \xrightarrow{\eps} L(\one_\C)
\]
exhibits its target as a unit object in $\D$ over $\one_{\cO}$. Indeed, since $L^{\otimes}$ is surjective on objects and preserves cocartesian maps over $\cO^{\otimes}$ that exist in $\C^{\otimes}$, it also preserves all operadically cocartesian maps that exist in $\C^{\otimes}$.
\end{enumerate}

\end{remarks}




The following proposition summarises the main facts we need about unit objects. It is obtained by combining~\cite{HA}*{Proposition 3.2.1.8, Lemma 3.2.1.10 \& Proposition 3.4.2.1}: 

\begin{proposition}
\label{proposition:units}%
Let $p\colon \D^\otimes \to \Com^{\otimes}$ be an $\infty$-operad.
Suppose that $u\colon \emptyset \to \one_{\D}$ exhibits $\one_{\D}$ as a unit object. Then the map 
$\mathbb{E}_0^{\otimes} \to \D^{\otimes}$ encoded by $u$ extends in an essentially unique manner to a commutative algebra object 
\[
\wtl{\one}_\D\colon \Com^{\otimes} \to \D^{\otimes}, \quad \<1\>  \mapsto \one_{\D},
\]
often referred to as trivial algebra object. In addition, this commutative algebra object satisfies the following properties:
\begin{enumerate}
\item
$\wtl{\one}_\D$ sends every map in $\Com^\otimes$ to a $p$-cocartesian map in $\D^{\otimes}$.
\item
With this algebra structure, $\one_{\D}$ is initial in $\CAlg(\D)$. 
\item
The forgetful functor $\Mod_{\one_{\D}}(\D)^{\otimes} \to \D^{\otimes}$ is an equivalence. In other words, any object in $\D$ is a $\one_\D$-module in any essentially unique manner.
\end{enumerate}
\end{proposition}



\begin{proposition}
\label{proposition:action-of-mon}%
In the situation of Proposition~\reffour{proposition:units}, if
\[
\mu \colon \MCom^\otimes \to \D^\otimes
\]
is an $\infty$-operad map encoding an algebra-module pair $(\wtl{\one}_{\D},d)$ in which the algebra object is trivial then 
$\mu$ sends every arrow in $\MCom^{\otimes}$ to an $p$-cocartesian arrow in $\D^{\otimes}$. %
\end{proposition}
\begin{proof}
Since $\mu|_{\Com^{\otimes}}$ encodes the trivial algebra object $\one_{\D}$, it sends the map from $\emptyset$ to the algebra colour in $\MCom$ to an operadically $p$-cocartesian arrow in $\D^{\otimes}$. As a consequence, $\mu$ sends every map that can be obtained by concatenating this one with objects in $\MCom^{\otimes}$ to a $p$-cocartesian arrow. By direct inspection, this includes all active arrows in $\MCom^{\otimes}$ which lie over (active) injective maps $\<n\> \to \<m\>$ in $\Com^{\otimes}$. It hence suffices to show that $\mu$ sends maps lying over active surjective maps $\<n\> \to \<m\>$ in $\Com^{\otimes}$ to $p$-cocartesian maps. But any such map in $\MCom^{\otimes}$ admits a section which is itself of the injective kind, and so the desired result follows from~\cite{HTT}*{Proposition 2.4.1.7}.
\end{proof}










Given a map of $\infty$-operads $q\colon \C^{\otimes} \to \cO^{\otimes}$ and a unit object $\one_{\cO} \in \cO$ with associated trivial algebra object $\tilde{\one}_{\cO}\colon \Com^{\otimes} \to \cO^{\otimes}$, let us denote by $\C^\otimes_{\mon/\one_\cO}$ the full suboperad of $\C^{\otimes} \times_{\cO^{\otimes}} \Com^{\otimes}$ spanned by the symmetric monoidal objects over $\one_\cO$.


\begin{proposition}
\label{proposition:main-point}%
Let $q\colon \C^{\otimes} \to \cO^{\otimes}$ be a map of $\infty$-operads such that $\cO^{\otimes}$ admits a unit object $\one_{\cO}\in \cO$ and $\C^{\otimes}$ admits a unit $\one_{\C}$ over $\one_{\cO}$. 
Suppose given a map $\mu\colon \MCom^{\otimes} \to \cO^{\otimes}$ encoding an algebra-module pair $(\wtl{\one}_{\cO},b)$ whose algebra is trivial, and form the commutative square
\[
\begin{tikzcd}
\M_b^{\otimes} \ar[r]\ar[d] & \C^{\otimes} \ar[d] \\
\MCom^{\otimes} \ar[r,"\mu"] & \cO^{\otimes}
\end{tikzcd}
\]
where $\M_b^{\otimes}$ is the full suboperad of $\MCom^{\otimes}\times_{\cO^{\otimes}} \C^{\otimes}$ spanned by the objects $(a, c)$ with $a \in \MCom$ the algebra colour and $c \in \C$ symmetric monoidal over $\one_{\cO}$ and all the objects of the form $(m,c)$ where $m \in \MCom$ is the module colour.
Then the left vertical map is a cocartesian fibrations of $\infty$-operads. In particular, $\C^{\otimes}_{\mon/\one_{\cO}}$ is a symmetric monoidal $\infty$-category and $\M_b^{\otimes}$ encodes an action of $\C^{\otimes}_{\mon/\one_{\cO}}$ on $\C_b$.
\end{proposition}

The proof of Proposition~\reffour{proposition:main-point} will make use of the following lemma:

\begin{lemma}
\label{lemma:sm-closed-under-tensor}%
Let $q\colon \C^{\otimes} \to \cO^{\otimes}$ be a map of $\infty$-operads.
Let $\bet\colon (x,y) \to x\otimes y$ be an operadically cocartesian lift of the map $\mu_2\colon \<2\> \to \<1\>$ such that $\alp=q(\bet)\colon (a,b) \to a\otimes b$ is an operadically cocartesian lift of $\mu_2$ in $\cO^{\otimes}$. Then the following holds:
\begin{enumerate}
\item
If $a$ and $b$ are symmetric monoidal then $a\otimes b$ is symmetric monoidal.
\item
If $a,b$ and $a\otimes b$ are symmetric monoidal and $x,y$ are symmetric monoidal over $a,b$, respectively, then $x\otimes y$ is symmetric monoidal over $a\otimes b$.
\end{enumerate}
\end{lemma}
\begin{proof}
The first claim can then be deduced from the second by applying it to the $\infty$-operad map $\cO^{\otimes} \to \Com^{\otimes}$.
We now prove the second claim. We need to check that operadically $p$-cocartesian lifts of the active map $\<2\> \to \<1\>$ exist for every source of the form $(z,w)$ with $w \in \C$, and that these lifts map to operadically $r$-cocartesian maps in $\cO$. For this, consider the commutative square
\[
\begin{tikzcd}
\<3\>\ar[r]\ar[d] & \<2\> \ar[d] \\
\<2\> \ar[r] & \<1\>
\end{tikzcd}
\]
of active arrows in $\Com^{\otimes}$ where the two maps $\<3\> \to \<2\>$ are the two distinct active maps sending $1$ to $1$ and $3$ to $2$. Since both $x$ and $y$ are symmetric monoidal over $a$ and $b$ respectively
we can lift the left vertical and bottom horizontal arrows to operadically $p$-cocartesian arrows 
\[
(x,y,w) \to (x,y \otimes w) \to (x \otimes y \otimes w)
\]
in $\C^{\otimes}$ whose image in $\cO^{\otimes}$ is operadically $r$-cocartesian, where as above we use the tensor notation suggestively. Let $\gam$ denote the composite of these two maps.
Since $\beta\colon (x,y) \to x\otimes y$ is operadically cocartesian its concatenation $(\beta,w)\colon (x,y,w) \to (x\otimes y,w)$ with $w$ is cocartesian and so we may then factor $\gam$ as a composite
\[
(x,y,w) \xrightarrow{(\beta,w)} (x\otimes y,w) \to (x \otimes y \otimes w) .
\]
Since $\gam$ and the first arrow are both operadically cocartesian it follows that the second arrow $(x\otimes y,w) \to (x \otimes y \otimes w)$ is an operadically cocartesian lift of $\<2\> \to \<1\>$ by the cancellation property for cocartesian arrows. Its image in $\cO$ is also an operadically cocartesian lift of $\<2\> \to \<1\>$ by the same cancellation property. %


\end{proof}

\begin{proof}[Proof of Proposition~\reffour{proposition:main-point}]
The morphisms in $\MCom^{\otimes}$ are generated under concatenation and composition by the unit $\emptyset \to a$, the multiplication map $(a,a) \to a$, and the action map $(a,m) \to m$. It will hence suffice to show that any of these three arrows admits operaically cocartesian lifts in $\M_b^{\otimes}$ for any lift of their domain. Since $\mu \colon \MCom^\otimes \to \D^\otimes$ sends every arrow to an operadically cocartesian arrow over $\Com$ by Proposition~\reffour{proposition:action-of-mon} the desired result now follows form Lemmas~\reffour{lemma:absolute-is-relative} and \reffour{lemma:sm-closed-under-tensor}.
\end{proof}



Proposition~\reffour{proposition:main-point} allows to define the notion of idempotent objects in the $\infty$-operadic context, which we will use in the subsequent section to define and study the notion of smashing localisations:
\begin{definition}
Let $\cO^{\otimes}$ be an $\infty$-operad with unit $\one_{\cO}$
and $q\colon \C^{\otimes} \to \cO^{\otimes}$ a map of $\infty$-operads with unit 
$\one_{\C}$
over $\one_\cO$. 
We will say that a map $\eps\colon \one_{\C} \to e$ in $\C$ exhibits $e$ as \defi{idempotent over $\cO$}, or \defi{$\cO$-idempotent}, if the following three conditions hold:
\begin{enumerate}
\item
The map $q(\eps)$ is an equivalence in $\cO$.
\item
The object $e$ is symmetric monoidal over $\one_{\cO}$. %
\item
The map
$\eps \otimes \id\colon e \to e \otimes e$ is an equivalence in the symmetric monoidal $\infty$-category $\C_{\mon/\one_{\C}}$.
\end{enumerate}
\end{definition}

\begin{remark}
\label{remark:action-of-idempotent}%
Any symmetric monoidal object $e$ in $\C$ over $\one_{\cO}$ determines an endo-functor $e \otimes (-)\colon \C_a \to \C_a$ for any $a \in \cO$ via the action of 
$\C_{\mon/\one_{\C}}$
on $\C_a$ of Proposition~\reffour{proposition:main-point}.
If $\eps\colon \one_{\C} \to e$ exhibits $e$ is idempotent over $\cO$ then the resulting endo-functor of $\C_a$ inherits a natural transformation from $\id_{\C_a}$ which satisfying the conditions of~\cite{HTT}*{Proposition 5.2.7.4}, and so $e \otimes (-)$ exhibits its image as a left Bousfield localisation of $\C_a$. The map $\epsilon\colon \one_{\C} \to e$ can then be recovered as the local approximation of $\one_{\C}$ in $\C_{\one_{\C}}$. Collections of left Bousfield localisations which arise in this manner are always $\cO$-multiplicative, and hence refine to $\cO$-multiplicative left Bousfield localisations of $\infty$-operads 
by Proposition~\reffour{proposition:multiplicative-localisation}. These special types of left Bousfield localisations are called \emph{smashing localisations}, which we discuss next.
\end{remark}












\subsection{Smashing localisations of \texorpdfstring{$\infty$-}{}operads}
\label{subsection:smashing-localisations}%


In this subsection we develop the notion of smashing localisations in the context of $\infty$-operads.
Let
\[
\begin{tikzcd}
\C^{\otimes} \ar[dr,"q"]\ar[ddr,"p"']\ar[rr,"L^{\otimes}"] && \ar[dl] \D^{\otimes}\\
& \cO^{\otimes} \ar[d,"r"] & \\
& \Com^{\otimes} &
\end{tikzcd}
\]
be an $\cO$-multplicative left Bousfield localisation of $\infty$-operads and assume that %
$\cO^{\otimes}$ admits a unit $\emptyset \to \one_{\cO}$ and 
$\C^{\otimes}$ admits a unit $\emptyset \to \one_{\C}$ over $\one_{\cO}$.

\begin{construction}
\label{construction:eta}%
Write $e \coloneq L(\one_{\C})$ and fix an object $x \in \C$ lying above a colour $a \in \cO$. 
Then the map $x \to (\one_\C,x)$ obtained by concatenating $\emptyset \to \one_{\C}$ with $x$ is $p$-cocartesian and maps to an $r$-cocartesian arrow in $\cO$. We may then factor the map $x \to L(x)$ as a composite in $\C^\otimes$ of the form appearing on the left
\[
\begin{tikzcd}[column sep = 10pt]
& (\one_\C,x)\ar[dr] &     &&  & (\one_{\cO},a)\ar[dr] &   &&  & \<2\> \ar[dr] & \\
x \ar[ur]\ar[rr] && L(x) &&  a \ar[ur]\ar[rr,equal] && a && \<1\> \ar[ur]\ar[rr,equal] && \<1\>
\end{tikzcd}
\]
lying over a triangle in $\cO^\otimes$ of the form appearing in the middle and over a triangle in $\Com^\otimes$ of the form appearing on the right, and where the arrow $(\one_{\cO},a) \to a$ is operadically $r$-cocartesian (by applying~\cite{HTT}*{Proposition 2.4.1.7} after concatenating with arbitrary objects in $\cO^{\otimes}$).
Since the map $\one_{\C} \to e$ is an $L$-equivalence, %
so is the map $(\one_{\C},x) \to (e,x)$, and so %
the map $(\one_{\C},x) \to L(x)$ factors through a map 
\[
\eta_x\colon (e,x) \to L(x)
\]
lying over the same operadically $r$-cocartesian arrow $(\one_{\cO},a) \to a$ (up to an equivalence). 
\end{construction}


\begin{definition}
\label{definition:smashing}%
Keeping the above notation and assumptions, we say that the $\cO$-multiplicative left Bousfield localisation $L^{\otimes}$ is \emph{$\cO$-smashing} if the map $\eta_x\colon (e,x) \to L(x)$ of Construction~\reffour{construction:eta} is operadically $p$-cocartesian for every $x \in \C$.
\end{definition}

The following is the main result of this subsection, which can be viewed as a recognition criterion as the one of Proposition~\reffour{proposition:multiplicative-localisation}, this time for smashing localisations.



\begin{proposition}
\label{proposition:equivalent-1}%
Let $\cO^{\otimes}$ be an $\infty$-operad with unit $\one_{\cO}$
and $q\colon \C^{\otimes} \to \cO^{\otimes}$ a $\infty$-operad over $\cO^{\otimes}$ which admits a unit $\emptyset \to \one_{\C}$ over $\one_\cO$. Let $\D^{\otimes} \subseteq \C^{\otimes}$ be a full suboperad.
Then the following are equivalent
\begin{enumerate}
\item
\label{item:smashing}%
The inclusion $\D^{\otimes} \subseteq \C^{\otimes}$ admits an $\cO$-smashing left Bousfield localisation $L^{\otimes}\colon \C^{\otimes} \to \D^{\otimes}$. %
\item
\label{item:idem-1}%
There exists an $\cO$-idempotent object $\eps\colon \one_{\C} \to e$ in $\C$ such that for every $a \in \cO$ the full subcategory $\D_a \subseteq \C_a$ is exactly the image of $e \otimes (-)\colon \C_a \to \C_a$. 
\item
\label{item:idem-2}%
Each inclusion $\D_a \subseteq \C_a$ admits a left adjoint $L_a\colon \C_a \to \D_a$
and there exists an $\cO$-idempotent object $\eps\colon\one_{\C} \to e$ in $\C$ such that $e \otimes (-) \simeq L_a(-)$ as endo-functors under $\id_{\C_a}$ for every $a \in \cO$. In particular, %
$e \simeq L_{\one_{\cO}}(\one_{\C})$ under $\one_{\C}$. %
\item
\label{item:closed}%
Each inclusion $\D_a \subseteq \C_a$ admits a left adjoint $L_a\colon \C_a \to \D_a$
such that the collection $\{L_a\}_{a \in \cO}$ is $\cO$-multiplicative, the object $e := L_{\one_{\cO}}(\one_{\C})$ is symmetric monoidal over $\one_{\cO}$,
and for any $a \in \cO$, $x \in \C_{\mon/\one_\cO}$ and $y \in \C_a$, if either $x$ belongs to $\D_{\one_{\cO}}$ or $y$ belongs to $\D_{a}$, then $x \otimes y \in \C_a$ belongs to $\D_a$. 
\end{enumerate}
\end{proposition}
\begin{proof}
We begin with proving \reffouritem{item:smashing} $\Rightarrow$ \reffouritem{item:idem-1}, and so we 
suppose that the inclusion $\D^{\otimes} \subseteq \C^{\otimes}$ admits an $\cO$-smashing left Bousfield localisation $L^{\otimes}\colon \C^{\otimes} \to \D^{\otimes}$.
Then by definition for every $x \in \C$ there exists an operadically cocartesian map $(e,x) \to L_{q(x)}(x)$ lifting $\<2\> \to \<1\>$ and which maps to an operadically cocartesian map in $\cO^{\otimes}$ over $\Com^{\otimes}$, so that $e$ is symmetric monoidal over $\one_{\cO}$
by Remark~\reffour{remarks:symmetric-monoidal-objects-and-units}\;\reffouritem{item:simplifies}. In addition, taking $x = e$ we obtain a triangle
\[
\begin{tikzcd}
& (e,e) \ar[dr] & \\
e \ar[rr,"\simeq"]\ar[ur] && L_{\one_{\cO}}(e)
\end{tikzcd}
\]
in which the right diagonal map is operadically cocartesian over $\Com^{\otimes}$ and the bottom horizontal map is an equivalence since $e$ belongs to $\D_{\one_{\cO}}$, and so the map $\one_{\C} \to e$ exhibits $e$ as idempotent over $\cO$. Finally, for every $a \in \cO$ and $x \in \C_a$ the $q$-cocartesian map $\eta_x \colon (e,x) \to L_a(x)$ induces an equivalence $e \otimes x \simeq L_a(x)$, and hence $\D_a$ coincides with the image of $e \otimes (-)\colon \C_a \to \C_a$, proving \reffouritem{item:idem-1}.
The implication \reffouritem{item:idem-1} $\Rightarrow$ \reffouritem{item:idem-2} follows from Remark~\reffour{remark:action-of-idempotent}. %
To prove \reffouritem{item:idem-2} $\Rightarrow$ \reffouritem{item:closed}, we assume that each inclusion $\D_a \subseteq \C_a$ admits a left adjoint $L_a\colon \C_a \to \D_a$ and that we have an $\cO$-idempotent $\one_{\C} \to e$ %
such that $\sig\colon e \otimes (-) \simeq L_a(-)$ as endo-functors under $\id_{\C_a}$ for every $a \in \cO$. Let us first show that $\{L_a\}_{a \in \cO}$ is $\cO$-multiplicative. For this, we need to show that for every operation $\alp\colon (a_1,...,a_n) \to b$ in $\cO$, any choice of an $L_{a_i}$-equivalence $f_i\colon x_i \to y_i$ in $\C_{a_i}$ for each $i=1,...,n$, and any $d \in \D_b$, the restriction map
\[
\Mul^{\alp}_{\C}(y_1,\dots,y_n;d) \to \Mul^{\alp}_{\C}(x_1,\dots,x_n; d)
\]
on spaces of operations in $\C^{\otimes}$ covering $\alp$, is an equivalence. By Remark~\reffour{remark:full-suboperad}~\reffouritem{item:one-at-a-time} we may assume without loss of generality that all the $f_i$ are identities (or just equivalences) except for $f_1$. 
Now consider the commutative diagram
\[
\begin{tikzcd}
\Mul^{\alp}_{\C}(y_1,\dots,y_n;d) \ar[r]\ar[d] & \Mul^{\alp}_{\C}(x_1,\dots,x_n; d) \ar[d] \\
\Mul^{\alp}_{\C}(e \otimes y_1,y_2,...,y_n;e \otimes d) \ar[r,"\simeq"]\ar[d] & \Mul^{\alp}_{\C}(e\otimes x_1,x_2,...,x_n; e\otimes d) \ar[d] \\
\Mul^{\alp}_{\C}(y_1,\dots,y_n;e \otimes d) \ar[r] & \Mul^{\alp}_{\C}(x_1,\dots,x_n; e\otimes d)  \\
\end{tikzcd}
\]
where the maps from the top row to the middle row are induced by tensoring with $e$ (well-defined since $e$ is a symmetric monoidal object over $\one_{\cO}$),
and the maps from the middle to the bottom row are induced by precomposition with $y_1 \to e \otimes y_1$ and $x_1 \to e \otimes x_1$, respectively. 
Then the middle horizontal map is an equivalence by the assumption that each $f_i\colon x_i \to y_i$ is an equivalence except for $f_1$, which is an $L_{a_1}$-equivalence (and given that we have an equivalence $L_{a_1}(-) \simeq e \otimes (-)$). At the same time, the composite vertical maps coincide with the ones induced by post-composition with $d \to e \otimes d$ by naturality of concatenation and of factorisation through cocartesian lifts, and is hence an equivalence by the assumption that $d$ belongs to $\D_a$. The entire diagram is hence equivalent to a retract diagram in the arrow category of spaces, and since the middle row is an equivalence we conclude that the top (and bottom) horizontal arrows are equivalences, as desired.

Moving forward, note that since $e \simeq L_{\one_{\cO}}(\one_\C)$ is assumed $\cO$-idempotent, it is in particular symmetric monoidal over $\one_{\cO}$.
To establish~\reffouritem{item:closed} it is hence left to show that for any $a \in \cO$, $x \in \C_{\mon/\one_\cO}$, and $y \in \C_a$, if either $x$ belongs $\D_{\one_{\cO}}$ or $y$ belongs to $\D_{a}$ then $x \otimes y \in \C_a$ belongs to $\D_a$. Since $L_a(-) \simeq e \otimes (-)$ this amounts to showing that the induced map 
\[
\eps\otimes x \otimes y\colon x \otimes y \to e \otimes x \otimes y
\]
is an equivalence. Now if $x$ belongs to $\D_{\one_{\cO}}$ then $\eps \otimes x \colon x \to e \otimes x$ is an equivalence and hence $\eps \otimes x \otimes y$ is an equivalence as well, and if $y$ belongs to $\D_a$ then $\eps \otimes y \colon y \to e \otimes y$ is an equivalence and hence $\eps \otimes x \otimes y \simeq x \otimes \eps \otimes y$ is an equivalence as well. 

Finally, we prove \reffouritem{item:closed} $\Rightarrow$ \reffouritem{item:smashing}. By Proposition~\reffour{proposition:multiplicative-localisation} we have that the collection $\{L_a\}_{a \in \cO}$ refines to a left Bousfield localisation $L^{\otimes}\colon \C^{\otimes} \to \D^{\otimes}$ of $\infty$-operads over $\cO^{\otimes}$. It is left to show that $L^{\otimes}$ is $\cO$-smashing. %
Let us set $e := L(\one_{\C})$ and for $a \in \cO$ and $x \in \C_a$ consider the commutative triangle
\[
\begin{tikzcd}
& (e,x) \ar[dr,"\eta_x"] & \\
x \ar[rr]\ar[ur] && L(x) 
\end{tikzcd}
\]
resulting from Construction~\reffour{construction:eta}. We need to show that $\eta_x$ is operadically cocartesian over $\Com^{\otimes}$. Let us first point out that since $e$ is assumed symmetric monoidal over $\one_{\cO}$
we may factor $\eta_x$ as a composite 
\[
(e,x) \to e \otimes x \xrightarrow{\eta_x'} L(x)
\]
where the first map is an operadically cocartesian lift of $\<2\> \to \<1\>$ and maps to an operadically cocartesian 
lift of this map in $\cO$, while $\eta_x'$ is a map that lies in the fibre $\C_a$ over $a \in \cO$. It will hence suffice to show that $\eta_x'$ is an equivalence. 
By assumption, we have that $e \otimes x$ belongs to $\D_a$. Since $L(x)$ also belongs to $\D_a$ %
it is enough to show that $\eta_x'$ is a $\D_a$-local equivalence, that is, that for every $d \in \D_a$ the induced map 
\[
\Map_{\C_a}(L(x),d) \to \Map_{\C_a}(e \otimes x,d) \simeq \Mul_{\C}(e,x; d)
\]
is an equivalence. By the multiplicativity of $\{L_a\}_{a \in \cO}$ and the unitality of $\one_{\C}$, the restriction maps
\[
\Mul_{\C}(e,x; d) \to \Mul_{\C}(\one_{\C},x; d) \to \Map_{\C_a}(x,d)
\]
are both equivalences. Finally, the total composite 
\[
\Map_{\C_a}(L(x),d) \to \Map_{\C_a}(x,d)
\]
coincides with restriction along $x \to L(x) = L_a(x)$, and is hence an equivalence for every $d \in \D_a$, completing the proof of the proposition.
\end{proof}





\begin{corollary}
\label{corollary:action-smashing}%
Let $\cO^{\otimes}$ be an $\infty$-operad with unit $\one_{\cO}$, $q\colon \C^{\otimes} \to \cO^{\otimes}$ a map of $\infty$-operads with unit $\one_{\C}$ over $\one_\cO$, and $L^{\otimes} \colon \C^{\otimes} \to \D^{\otimes}$ an $\cO$-smashing left Bousfield localisation. 
Write $e := L(\one_{\C})$ for the image of $\one_{\C}$, so that $e$ is a unit object in $\D$ by Remark~\reffour{remarks:symmetric-monoidal-objects-and-units}~\reffouritem{item:local-unit}.
Then in the square
\[
\begin{tikzcd}
\Mod_e(\D)^{\otimes} \ar[r]\ar[d,"\simeq"'] & \Mod_e(\C)^{\otimes} \ar[d]\\
\D^{\otimes} \ar[r] & \C^{\otimes},
\end{tikzcd}
\]
where the left vertical map is an equivalence by Proposition~\reffour{proposition:units},
the top horizontal map is also an equivalence, yielding an equivalence $\D^{\otimes} \simeq \Mod_e(\C)^{\otimes}$ of $\infty$-operads over $\C^{\otimes}$.
\end{corollary}
\begin{proof}
Since $\D^{\otimes} \to \C^{\otimes}$ is a full suboperad inclusion, the top horizontal map in the displayed diagram is again a full suboperad inclusion. It will hence suffice to show that it is essentially surjective. In other words, we need to show that if $x \in \C^{\otimes}$ refines to an $e$-module object then $x$ lies in $\D$. Now, by Proposition~\reffour{proposition:equivalent-1}, the map $\eps\colon \one_{\C} \to e$ exhibits $e$ as idempotent over $\cO$, and so $e$ is in particular symmetric monoidal over $\one_{\cO}$.
We may consequently view $e$ as a commutative algebra object in the symmetric monoidal $\infty$-category $\C_{\mon/\one_\cO}$. 
Since $e$ lies over the trivial algebra object $\one_{\cO}$ in $\cO^{\otimes}$, the object $a$ admits an essentially unique $\one_{\cO}$-module structure, and so any $e$-module structure on $x$ will induce this structure on $a := q(x)$ upon applying $q$. 
In particular, any $e$-module structure on $x$ uniquely refines to an $e$-module structure on $x$ as an object of $\C_a$, where we view $\C_a$ as a module over the symmetric monoidal $\infty$-category $\C_{\mon/\one_\cO}$ (see Proposition~\reffour{proposition:action-of-mon}) and $e$ as a commutative algebra object there. In particular, if $x$ is an $e$-module in $\C_a$ then the unit and action maps
\[
x \to e \otimes x \to x
\]
exhibit $x$ as a retract of $e \otimes x$. Applying Proposition~\reffour{proposition:equivalent-1} again we have that $e \otimes x$ belongs to $\D_a$, and hence $x$ belongs to $\D_a$ as well (since reflective subcategories are closed under retracts).  
\end{proof}


Finally, let us also summarise the conclusions of the above discussion in the case where $\cO^{\otimes} = \Com^\otimes$ and $\C^{\otimes}$ is a symmetric monoidal $\infty$-category.

\begin{corollary}
\label{corollary:smashing-symmetric-monoidal}%
Let $L^{\otimes} \colon \C^{\otimes} \to \D^{\otimes}$ be a left Bousfield localisation of symmetric monoidal $\infty$-categories with fully faithful (and lax symmetric monoidal) right adjoint $\iota^{\otimes}\colon \D^{\otimes} \hrar \C^{\otimes}$.
Then the following conditions are equivalent:
\begin{enumerate}
\item
\label{item:counit-tensor}%
$L$ is smashing. %
\item
\label{item:left-adjoint-is-tensoring}%
There exists an idempotent $e \in \C_{\one/}$ such that $\iota L$ is equivalent, as a functor under $\id_{\C}$, to $x \mapsto e \otimes x$.
\item
\label{item:closed-under-tensoring}%
The image of $\iota\colon \D \hrar \C$ is closed under tensoring with arbitrary objects in $\C$, that is, it is an ideal.
\item
\label{item:smashing-as-algebra}%
There exists an idempotent algebra object $e \in \CAlg(\C^{\otimes})$ and a symmetric monoidal equivalence $\D^{\otimes} \simeq \Mod_e(\C)^{\otimes}$ under which $\iota$ identifies with the forgetful functor $\Mod_e(\C) \to \C$.
\end{enumerate}
\end{corollary}
\begin{proof}
The equivalence of the first three condition follows from Proposition~\reffour{proposition:equivalent-1}, while Corollary~\reffour{corollary:action-smashing} shows that~\reffouritem{item:smashing-as-algebra} follows from~\reffouritem{item:counit-tensor}. To finish the proof it will hence suffices to note that~\reffouritem{item:smashing-as-algebra} implies~\reffouritem{item:closed-under-tensoring}, since the forgetful functor $\Mod_e(\C) \to \C$ refines to one of $\C$-modules. %
\end{proof}



\begin{remark}
\label{remark:smashing-closed-colimits}%
In the situation of Corollary~\reffour{corollary:smashing-symmetric-monoidal}, if $\C$ admits small colimits and its tensor product preserves small colimits in each variable then $\iota(\D) \subseteq \C$ is closed under colimits by the characterisation in \reffouritem{item:smashing-as-algebra}. Let us emphasise that the latter condition is strictly weaker than being smashing, see Example~\reffour{example:cant-remove} below for a multiplicative localisation closed under colimits which is not smashing. Nonetheless, since in any multiplicative localisation $\iota(\D)$ is closed under tensoring with dualisable objects, if $\C$ is generated under colimits by dualisable objects and $\iota(\D) \subseteq \C$ is closed under colimits then the localisation is smashing by the characterisation in~\reffouritem{item:closed-under-tensoring}. This is the case, for example, when $\C=\Spa$, whose smashing localisations have been the subject of much study.
\end{remark}





\subsection{Multiplicative localisations of Day convolution}
\label{subsection:day-convolution}%
\label{subsec:large-indexing}%

In this subsection we specialise the general discussion above to localisations of Day convolution $\infty$-operads obtained by requiring a certain collection of cone diagrams to be sent to limits. Our main goal is to establish Propositions~\reffour{proposition:multiplicative-day} and \reffour{proposition:smashing-day} below, which give sufficient criterion for such localisations to be multiplicative and smashing, respectively. 

  %

Let $\cO^{\otimes}$ be an $\infty$-operad and $\C^{\otimes} \to \cO^{\otimes}$ an $\cO$-monoidal $\infty$-category. Then 
the functor
\[
(-) \times_{\cO^{\otimes}} \C^{\otimes} \colon (\Op)_{/\cO} \to (\Op)_{/\cO}
\]
admits a right adjoint
\[
\Fun_{\cO}(\C,-)^{\Day}\colon (\Op)_{/\cO} \to (\Op)_{/\cO}
\]
given by the formation of Day convolution, see~\cite{HA}*{Construction 2.2.6.7}; an alternative proof is given in~\cite{hinich-enriched}*{Proposition 2.8.3}, which also gives a simple explicit model in the case where both $\C^{\otimes}$ and $\D^{\otimes}$ are $\cO$-monoidal. The latter context is in fact sufficient for our purposes, and we shall henceforth fix the assumption that $\D^{\otimes}$ (in addition to $\C^{\otimes}$) is $\cO$-monoidal.
The universal property of Day convolution then gives %
\[
\Alg_{\cP/\cO}(\Fun_{\cO}(\C,\D)^{\Day}) = \Alg_{\cP \times_{\cO} \C/\cO}(\D) =  \Alg_{\cP \times_{\cO} \C/\cP}(\cP \times_{\cO}\D) ,
\]
that is, the $\infty$-category of $\cP$-algebras in $\Fun_{\cO}(\C,\D)^{\Day}$ is equivalent to the $\infty$-category of lax $\cP$-monoidal functors $\C \to \D$ for any $\infty$-operad $\cP^{\otimes} \to \cO^{\otimes}$ over $\cO^{\otimes}$. 


For every $a \in \cO$, the $\infty$-category of colours of $\Fun_{\cO}(\C,\D)^{\Day}$ lying over $a$ is given by the functor category $\Fun(\C_a,\D_a)$. The spaces of operations in $\Fun_{\cO}(\C,\D)^{\Day}$ can then be described explicitly as follows. Given an operation $\alp\colon (a_1,...,a_n) \to b$ in $\cO$ with associated $\cO$-monoidal operation $\alp_{\D}\colon \D_{a_1} \times \dots \times \D_{a_n} \to \D_b$, functors $f_i\colon \C_{a_i} \to \D_{a_i}$ for $i=1,...,n$ and a functor $g \colon \C_b \to \D_b$, the corresponding space of operations over $\alp$ is given by the space of natural transformations
\[
\Mul^{\alp}_{\Fun_{\cO}(\C,\D)}(f_1,...,f_n;g) = \Nat(f_1 \boxtimes \dots \boxtimes f_n,g|_{\C_{a_1} \otimes \dots \otimes \C_{a_n}} ),
\]
between the composite 
\[
f_1 \boxtimes \dots \boxtimes f_n\colon \C_{a_1} \times \dots \times \C_{a_n} \to \D_{a_1} \times \dots \times \D_{a_n} \xrightarrow{\alp_{\D}} \D_b
\]
and the restriction of $g$ along the $\alp$-product
$\alp_{\C}\colon \C_{a_1} \times \dots \times \C_{a_n} \to \C_{b}$ in $\C^{\otimes}$. %



\begin{definition}
\label{definition:cocomplete}%
Let $\D^{\otimes} \to \cO^{\otimes}$ be an $\cO$-monoidal $\infty$-category.  
\begin{enumerate}
\item
We say that $\D^{\otimes}$ is \defi{$\cO$-cocomplete} if $\D_a$ admits small colimits for every $a \in \cO$ and $\alp_{\D}\colon \D_{a_1} \times ... \times \D_{a_n} \to \D_b$ is entry-wise colimit preserving 
for every arrow $\alp\colon (a_1,...,a_n) \to b$ in $\cO^{\otimes}$ with target in $\cO$.
\item
We say that $\D^{\otimes}$ is \defi{$\cO$-closed} if for every $\alp\colon (a_1,...,a_n) \to b$ in $\cO^{\otimes}$ with target in $\cO$ and every $(d_2,\dots,d_n) \in \D_{a_2} \times \dots \times \D_{a_n}$, the operation $\alp_{\D}(-,d_2,...,d_n)\colon \D_{a_1} \to \D_b$ admits a right adjoint $R^{d_2,...,d_n}_{\alp}\colon \D_b \to \D_{a_1}$;
\item We say that $\D$ is \defi{presentably} $\cO$-monoidal if it is $\cO$-cocomplete and each $\D_a$ is presentable.
\end{enumerate} 
\end{definition}

By the adjoint functor theorem, any presentably $\cO$-monoidal $\infty$-category is $\cO$-closed. 





The following proposition is a variant of~\cite{HA}*{Proposition 2.2.6.16}. We include the proof for the convenience of the reader.

\begin{proposition}
\label{proposition:day-cocartesian}%
Let $\D^{\otimes} \to \cO^{\otimes}$ be an $\cO$-cocomplete $\cO$-monoidal $\infty$-category and $\C$ a small $\cO$-monoidal $\infty$-category. Then the Day convolution $\infty$-operad $\Fun_{\cO}(\C,\D)^{\Day} \to \cO^{\otimes}$ is an $\cO$-monoidal $\infty$-category.
\end{proposition}

In fact, in the situation of Proposition~\reffour{proposition:day-cocartesian} one can show that $\Fun_{\cO}(\C,\D)^{\Day}$ is again $\cO$-cocomplete, and is further presentably $\cO$-monoidal if $\D$ is so. %

\begin{proof}[Proof of Proposition~\reffour{proposition:day-cocartesian}]
An arrow $\alp\colon (a_1,...,a_n) \to b$ in $\cO^{\otimes}$ with target in $\cO$ admits a locally cocartesian lift in $\Fun_{\cO}(\C,\D)^{\Day}$ with source $(f_1,...,f_n)$ if and only if 
\[
f_1 \boxtimes \dots \boxtimes f_n\colon \C_{a_1} \times ... \times \C_{a_n} \to \D_b
\]
admits a left Kan extension along $\alp_{\C}\colon \C_{a_1} \times ... \times \C_{a_n} \to \C_b$. This holds when $\C$ is small and each $\D_a$ admits small colimits. 
Theses locally cocartesian lifts then remain locally cocartesian after concatenation with fixed objects (since locally cocartesian arrows are closed under concatenation in any $\infty$-operad). The condition that all $\cO$-monoidal operation functors preserves small colimits in each entry then implies that all these locally cocartesian arrows are actually cocartesian.
Since cocartesian arrows are closed under composition we get that $\Fun_{\cO}(\C,\D)^{\Day} \to \cO^{\otimes} \to \cO^{\otimes}$ admits a sufficient supply of cocartesian lifts, as desired.
\end{proof}


\begin{examples}\
\begin{enumerate}
\item
Suppose $\cO=\Com$, $\C$ is a small symmetric monoidal $\infty$-category and $\D$ is a $\Com$-cocomplete symmetric monoidal $\infty$-category, that is, $\D$ admits small colimits and the monoidal product preserves small colimits in each entry. 
Then $\Fun_{\Com}(\C,\D)^{\Day}$ is the symmetric monoidal $\infty$-category whose underlying $\infty$-category is the $\infty$-category of functors $\C \to \D$ and whose tensor product $(f,g) \mapsto f \otimes_{\Day} g$ is given by a left Kan extension of the form 
\[
\begin{tikzcd}
\C \times \C \ar[d]\ar[r, "f \times g"] & \D \times \D \ar[r, "\otimes"] & \D \\
\C \ar[urr, dotted] && \ .
\end{tikzcd}
\]
This product is what is classically known as \emph{Day convolution}, see~\cite{day1970} for the original treatment of the idea in the setting of ordinary categories. 
symmetric monoidal $\infty$-category.
\item
Suppose that $\cO = \MCom$ (see Example~\reffour{example:mcom-loc}), $\C = (\C_a,\C_m)$ is small $\MCom$-monoidal $\infty$-category and $\D = (\cB,\E)$ is an $\MCom$-monoidal $\infty$-category %
such that $\cB$ is a $\Com$-cocomplete
symmetric monoidal $\infty$-category and $\E$ is a $\cB$-module in cocomplete $\infty$-categories.
Then the $\MCom$-monoidal $\infty$-category $\Fun^{\MCom}(\C,\D)^{\Day}$ encodes the data of the pair $(\Fun(\C_a,\cB),\Fun(\C_m,\E))$ where $\Fun(\C_a,\cB)$ is the %
symmetric monoidal functor category with its Day convolution as above, and $\Fun(\C_m,\E)$ is a $\Fun(\C_a,\cB)$-module %
with action of $f$ on $g$ given similarly by a left Kan extension of the form 
\[
\begin{tikzcd}
\C_a \times \C_m \ar[d]\ar[r, "f \times g"] & \cB \times \E \ar[r, "\otimes"] & \E \\
\C_m \ar[urr, dotted] && \ .
\end{tikzcd}
\]
\end{enumerate}
\end{examples}



\begin{remark}
\label{remark:left-kan-monoidal}%
Suppose given an $\cO$-monoidal functor $\rho^{\otimes}\colon \cB^{\otimes} \to \C^{\otimes}$ among small $\cO$-monoidal $\infty$-categories and let $\D^{\otimes} \to \cO^{\otimes}$ be an $\cO$-cocomplete $\cO$-monoidal $\infty$-category, so that $\Fun_{\cO}(\cB,\D)^{\Day}$ and $\Fun_{\cO}(\C,\D)^{\Day}$ are $\cO$-monoidal by Proposition~\reffour{proposition:day-cocartesian}. Then $\rho$ induces a (solidly indicated) restriction functor 
\[
\begin{tikzcd}
\Fun_{\cO}(\cB,\D)^{\Day} \ar[rr,bend left =15,dotted,"\rho_!"]\ar[dr]&& \ar[ll,"\rho^*","\perp"'] \Fun_{\cO}(\C,\D)^{\Day}\ar[dl] \\
& \cO^{\otimes} &
\end{tikzcd}
\]
which is generally only a map of $\infty$-operads over $\cO$ (that is, a lax $\cO$-monoidal functor), where for every $a \in \cO$, the induced map $\Fun(\C_a,\D_a) \to \Fun(\cB_a,\D_a)$ is given by restriction along $\rho$. Since each $\D_a$ is assumed cocomplete each of these restriction functors admits a left adjoint given by left Kan extension,  
and these left Kan extensions assemble to form a relative left adjoint $\rho_!$ over $\cO$ by~\cite{HA}*{Corollary 7.3.2.12}, which is automatically %
an $\cO$-monoidal functor (since relative left adjoints preserve relative colimits). 
We may summarise the situation as follows: Left Kan extension along symmetric monoidal functors among small symmetric monoidal $\infty$-categories induces an $\cO$-monoidal functor on the level of Day convolution $\infty$-operads
with $\cO$-cocomplete targets. 
\end{remark}






To start the passage to arbitrary locally small
source category $\C$, note that if $\C_0 \subseteq \C$ is some small full subcategory and $f\colon \C_0 \to \D$ is a functor with $\D$ cocomplete, then $f$ admits a left Kan extension to all of $\C$ by the pointwise formulae. We refer to functors which are obtained in the manner as \defi{small}, and write $\Fun^{\sm}(\C,\D) \subseteq \Fun(\C,\D)$ for the full subcategory spanned by the small functors. These are all generated under small colimits from functors of the following form:

\begin{notation}
\label{notation:representable}%
Given $c \in \C$ and $d \in \D$, we will denote by $c_!d \colon \C \to \D$ the functor obtained by left Kan extending the functor $\{c\} \to \D$ sending $c$ to $d \in \D$ along the inclusion $\{c\} \subseteq \C$. 
\end{notation}

Explicitly, $c_!d$ is given on some object $c' \in \C$ by the formula $d_{\Hom_\C(c,c')} \simeq \colim_{\Hom_\C(c,c')} \const_d$.

\begin{lemma}
\label{lemma:canonical-colimit}%
Let $f\colon \C \to \D$ be a functor from a locally small $\infty$-category $\C$ to an $\infty$-category with colimits over (small) spaces, and let $\vphi\colon \Twar(\C) \to \Fun(\C,\D)$ be the functor given by the formula 
\[
[x \to y] \mapsto y_!f(x) .
\]
Then the canonical cone diagram $\ovl{\vphi}\colon \Twar(\C)^{\triangleright} \to \Fun(\C,\D)$ on $\vphi$ with cone point $f$ is a colimit diagram.
\end{lemma}

We note that in the situation of Lemma~\reffour{lemma:canonical-colimit} the $\infty$-category $\Twar(\C)$ is potentially not small, and hence the colimit of $\vphi$ does not exist a-priori even if we ask $\D$ to be cocomplete, only a-posteriori. 

\begin{remark}
\label{remark:canonical-colimit}%
If $f$ is left Kan extended from some small full subcategory $\C_0 \subseteq \C$ then we may apply this lemma to $f|_{\C_0}$ to deduce that $f$ is a small colimit of functors of the form $c_!d$.
\end{remark}

\begin{proof}[Proof of Lemma~\reffour{lemma:canonical-colimit}]
Since colimits in functor $\infty$-categories that exists pointwise are also computed pointwise it will suffice to show that for every $z \in \C$ the corresponding cone diagram
\[
\ev_{z} \circ \ovl{\vphi}\colon \Twar(\C)^{\triangleright} \to \D
\]
with cone point $f(z)$ is a colimit diagram. Explicitly, this diagram sends $[x \to y] \in \Twar(\C)$ to $f(x)_{\Hom_{\C}(y,z)} \in \D$. In particular, we see that it is the left Kan extension along the projection 
\[
\Twar(\C_{/z}) \simeq \Twar(\C) \times_{\C\op} (\C_{/z})\op \to \Twar(\C)
\]
of the functor $[x \to y \to z] \mapsto f(x)$. It will hence suffice to show that the corresponding cone diagram
\[
\Twar(\C_{/z})^{\triangleright} \to \D
\]
with cone point $f(z)$ is a colimit diagram. We now observe that both the functor $[x \to y \to z] \mapsto f(x)$ and the natural transformation from it to $f(z)$ do not depend on $y$, and so the last cone diagram factors as
\[
\Twar(\C_{/z})^{\triangleright} \to \C_{/z}^{\triangleright} \to \D
\]
where the first functor sends $[x \to y \to z]$ to $[x \to z]$. Now $\Twar(\C_{/z}) \to \C_{/z}$ is a cartesian fibration with contractible fibres and is hence cofinal, see e.g.\ \cite{HHLN-two-var-fil}*{Lemma 5.5}. It will hence suffice to show that the cone diagram
\[
\C_{/z}^{\triangleright} \to \D
\]
is a colimit diagram. Indeed, this is because $\C_{/z}$ has a final object and the above diagram sends the arrow from the final object to the cone point to the identity on $f(z)$. 
\end{proof}



Now unlike the situation in Proposition~\reffour{proposition:day-cocartesian}, when $\C$ is only locally small the Day $\infty$-operad $\Fun_{\cO}(\C,\D)^{\Day}$ is generally not an $\cO$-monoidal $\infty$-category, even if $\D$ is $\cO$-cocomplete. In the following proposition and subsequent corollary we summarise a few properties which nonetheless persist to the locally small setting:


\begin{proposition}
\label{proposition:persist-locally-small}%
Let $\C$ be a locally small $\cO$-monoidal $\infty$-category and $\D$ an $\cO$-cocomplete $\cO$-monoidal $\infty$-category. 
Let $\alp\colon (a_1,\dots,a_n) \to b$ be an operation in $\cO$ and $f_i\colon \C_{a_i} \to \D_{a_i}$ be functors. Then $\alp$ admits an operadically cocartesian lift to $\Fun_{\cO}(\C,\D)$ starting from $(f_1,...,f_n)$ in either of the following cases:
\begin{enumerate}
\item
\label{item:all-small}%
The $f_i$ are all small. In particular, if $n=0$.
\item
The $f_i$ are small for all $i$ except at most one $i_0 \in \{1,...,n\}$ and $\C^{\otimes}$ is $\cO$-closed.
\end{enumerate}
In addition, in case~\reffouritem{item:all-small} the target of the resulting operadically cocartesian lift is again small.
\end{proposition}

\begin{corollary}
\label{corollary:persist-locally-small}%
Let $\C$ be a locally small $\cO$-monoidal $\infty$-category and $\D$ an $\cO$-cocomplete $\cO$-monoidal $\infty$-category. Then the following holds:
\begin{enumerate}
\item
The full suboperad $\Fun_{\cO}^{\sm}(\C,\D)^{\Day} \subseteq \Fun_{\cO}(\C,\D)^{\Day}$ spanned by $\Fun^{\sm}(\C,\D_a) \subseteq \Fun(\C,\D_a)$ for $a \in \cO$ is an $\cO$-monoidal $\infty$-category.
\item
If $\cO$ admits a unit object $\one_{\cO}$ then $\Fun_{\cO}(\C,\D)^{\Day}$ admits a unit object lying above $\one_{\cO}$.
\item
\label{item:small-monoidal}%
If $\C$ is $\cO$-closed and $a \in \cO$ is a symmetric monoidal object then any small $f\colon \C_a \to \D_a$ is a symmetric monoidal object in $\Fun_{\cO}(\C,\D)^{\Day}$ lying over $a$.
\end{enumerate}
\end{corollary}


\begin{proof}[Proof of Proposition~\reffour{proposition:persist-locally-small}]
Showing the existence of the desired operadically cocartesian lift amounts to showing that there exists a dashed left Kan extension 
\[
\begin{tikzcd}[column sep = 50pt]
\C_{a_1} \times \dots \times \C_{a_n} \ar[r,"f_1 \times ... \times f_n"]\ar[d,"\alp_{\C}"'] & \D_{a_1} \times \dots \times \D_{a_n} \ar[d,"\alp_{\D}"] \\
\C_b \ar[r,dashed] & \D_b
\end{tikzcd}
\]
of the top right composite along the left vertical arrow, which furthermore remains a left Kan extension after composing with all functors of the form $\beta(-,d_1,...,d_r)\colon \D_b \to \D_{b'}$ for operations $\beta\colon (b,b_1,...,b_r) \to b'$ and objects $d_i \in \D_{b_i}$. 
Let now $i_0 \in \{1,...,n\}$ be such that $f_i$ is small for $i \neq i_0$ and that either $f_{i_0}$ is also small or $\C$ is $\cO$-closed. The assumption that $\D$ is $\cO$-cocomplete implies that $\D_{a_i}$ and $\D_b$ admit small colimits and $\alp_D$ %
preserves small colimits in each argument, so that %
the top right composite depends on $f_1,...,f_n$ in a manner that preserves small colimits in each argument. Since $\D_b$ admits small colimits and the functors of the form $\beta(-,d_1,...,d_r)$ preserve small colimits by assumption the collection of $(f_1,....,f_n)$ for which the desired statement hold is closed under small colimits in each $f_i$ separately. Now by Lemma~\reffour{lemma:canonical-colimit} and Remark~\reffour{remark:canonical-colimit} small functors are all obtained from the basic ones of Notation~\reffour{notation:representable} via small colimits, and so we may assume without loss of generality that $f_{i} = (c_i)_!d_i$ for some $c_i \in \C,d_i \in \D$ for every $i \neq i_0$.
Using again that $\alp_{\D}$ commutes with small colimits in each argument we then get that 
\[
(c_1)_!d_1 \boxtimes \dots \boxtimes f_{i_0} \boxtimes \dots \boxtimes (c_n)_!d_n \colon \C \times \dots \times \C \to \D_{a_1} \times \dots \times \D_{a_n} \xrightarrow{\alp_{\D}} \D_b
\]
is left Kan extended from $\{c_1\} \times \dots \times \C \times \dots \times \{c_n\}$, and the same hold after post-composing with functors of the form $\beta(-,d_1,...,d_r)\colon \D_b \to \D_{b'}$.
It will hence suffice to show that the left Kan extension of $\alp_{\D}(d_1,...,f_{i_0},...,d_n)\colon \C_{a_{i_0}} \to \D_b$
along the functor $\alp_{\C}(c_1,...,-,...,c_n)\colon \C_{a_{i_0}} \to \C_b$
exists and preserved by post-composing with all functors of the form $\beta(-,d_1,...,d_r)$. If $\C$ is $\cO$-closed then the functor $\alp_{\C}(c_1,...,-,...,c_n)$ admits a right adjoint $R_{\alp}^{c_1,...,c_n}\colon \C_b \to \C_{a_{i_0}}$, and so left Kan extension along $\alp_{\C}(c_1,...,-,...,c_n)$ is given by pre-composition with $R_{\alp}^{c_1,...,c_n}$, an operation which commutes with post-composition with any functor. If $\C$ is not $\cO$-closed, then by assumption $f_{i_0}$ is also small, and so the argument above could also be applied at $i_0$ to reduced to the case where $f_{i_0} = (c_{i_0})_!d_{i_0}$ for some $c_{i_0} \in \C_{a_{i_0}}$ and $d_{i_0} \in \D_{a_{i_0}}$, in which case it will suffice to show that the functor $\ast \to \D_b$ picking the object 
$\alp_{\D}(d_1,...,d_n) \in \D_b$ admits a left Kan extension to $\C_b$ along the functor $\ast \to \C_b$ picking the object $\alp_{\C}(c_1,...,c_n)$, and that this left Kan extension is preserved by post-composition with any $\beta(-,d_1,...,d_r)$. Indeed, since $\C_b$ is locally small such left Kan extension exists as a pointwise one with small comma $\infty$-categories, and so the desired result follows. In addition, in this case we also get that the desired left Kan extension is itself small by definition, so that the last statement of Proposition~\reffour{proposition:persist-locally-small} follows from the fact that the collection of small functors is closed under small colimits.
\end{proof}


\begin{remark}
\label{remark:units}%
Unwinding the proof of Proposition~\reffour{proposition:persist-locally-small}
shows that if $\alp\colon \emptyset\to b$ is some nullary operation in $\cO$ with corresponding unit objects $\one^{\alp}_{\C} \in \C_b$ and $\one^{\alp}_{\D} \in \D_b$, then the corresponding unit object in $\Fun(\C_b,\D_b)$ is given by $(\one^{\alp}_{\C})_!\one^{\alp}_{\D}$ using Notation~\reffour{notation:representable}. In particular, if $\alp$ exhibits $b$ as a unit object of $\cO$ then the induced map $\emptyset \to (\one^{\alp}_{\C})_!\one^{\alp}_{\D}$ exhibits its target as a unit object in $\Fun_{\cO}(\C,\D)^{\Day}$ over $b$,
see Definition~\reffour{definition:monoidal-and-unit}. In this case, by~\cite{HA}*{Proposition 3.2.1.8 and Lemma 3.2.1.10} the objects $b$ and $(\one^{\alp}_{\C})_!\one^{\alp}_{\D}$ carry essentially unique commutative algebra structures, inducing symmetric monoidal structures on $\C_b$ and $\D_{b}$ together with a compatible lax symmetric monoidal structure on the functor $(\one^{\alp}_{\C})_!\one^{\alp}_{\D} \colon \C_b \to \D_b$ (by the universal property of Day convolution).
\end{remark}









\begin{definition}
\label{definition:almost-presentable}%
\label{definition:presentable}%
Let $\D^{\otimes} \to \cO^{\otimes}$ be an $\cO$-monoidal $\infty$-category. We say that $\D$ is \defi{almost closed} if it is $\cO$-cocomplete and every $\D_a$ is generated under small colimits by a 
(not necessarily small) full subcategory $\D'_a \subseteq \D_a$ 
satisfying the following property: for every operation $\alp\colon (a_1,...,a_n) \to b$ in $\cO$ and every $(d_2,\dots,d_n) \in \D'_{a_2}\times \dots \times \D'_{a_n}$, the operation $\alp_{\D}(-,d_2,...,d_n)\colon \D_{a_1} \to \D_b$ admits a right adjoint $R^{d_2,...,d_n}_{\alp}\colon \D_b \to \D_{a_1}$.
\end{definition}

\begin{remark}
Every closed $\cO$-monoidal $\infty$-category is almost closed. In particular, every presentably $\cO$-monoidal $\infty$-category is almost closed.
\end{remark}


\begin{example}
\label{example:mcom-almost-presentable}%
Suppose that $\cO^{\otimes} = \MCom^{\otimes}$, so that the data of a cocartesian fibration $\D^{\otimes} \to \MCom^{\otimes}$ corresponds to a pair $(\C,\M)$ where $\C$ is a symmetric monoidal $\infty$-category and $\M$ is a $\C$-module. Now if $\M$ is any cocomplete $\infty$-category, then $\M$ is canonically a module over $\Sps$ in a small colimit compatible manner (since $\Sps$ is the unit for the Lurie tensor product on cocomplete $\infty$-categories) in which case the pair $(\Sps,\M)$ is an $\MCom$-cocomplete $\MCom$-monoidal $\infty$-category. We claim that such an $(\Sps,\M)$ is almost closed. 
To see this, note that for every $m \in \M$ the functor $(-) \otimes m\colon \Sps \to \M$ admits a right adjoint $\Map_\M(m,-)$, and $\Sps$ is generated under colimits by its unit $\ast \in \Sps$, whose action $\ast \otimes (-) = \id_\M$ admits a right adjoint. Similarly, if $\M$ is a cocomplete additive $\infty$-category then $\M$ is canonically a module over $\Grp_{\Einf}$ in a manner compatible with small colimits and the pair $(\Grp_{\Einf},\M)$ constitutes an almost closed %
$\MCom$-monoidal $\infty$-category,
and if $\M$ is a cocomplete stable $\infty$-category then $\M$ is canonically a module over $\Spa$ in a manner compatible with small colimits and the pair $(\Spa,\M)$ constitutes an almost closed 
$\MCom$-monoidal $\infty$-category.
\end{example}



\begin{notation}
\label{notation:cone-diagrams}%
Let $\C$ be a locally small $\infty$-category equipped with a collection of cone diagrams $\cR = \{p_{\alp}\colon K_{\alp}^{\triangleleft} \to \C\}_{\alp}$. For any $\infty$-category %
$\E$ we write
$\Fun^{\cR}(\C,\E) \subseteq \Fun(\C,\E)$ for the full subcategory spanned by the $\cR$-local functors, that is, those functors $\C \to \E$ which send every cone diagram in $\cR$ to a limit cone. We will say that a natural transformation $\vphi \Rightarrow \vphi'$ is an $\cR$-local equivalence if it induces an equivalence on spaces of natural transformations to any $\cR$-local functor $\psi$. We will say that $\vphi \Rightarrow \vphi'$ is an $\cR$-local approximation if it is an $\cR$-local equivalence and $\vphi'$ is $\cR$-local.
We will say that a functor $\vphi\colon \C \to \E$ is \defi{$\cR$-localisable}  if it admits an $\cR$-local approximation $\vphi \Rightarrow \vphi'$. 
In particular, any $\cR$-local functor is $\cR$-localisable with the identity as an $\cR$-local approximation. We then write $\Fun^{\cR-\mathrm{able}}(\C,\E) \subseteq \Fun(\C,\E)$ for the full subcategory spanned by the $\cR$-localisable functors. We note that the inclusion $\Fun^{\cR}(\C,\E) \subseteq \Fun^{\cR-\mathrm{able}}(\C,\E)$ is reflective by design.

Now let $\cO$ be an $\infty$-operad, $\C^{\otimes},\D^{\otimes}$ two $\cO$-monoidal $\infty$-categories and let $\cR = \{\cR_a\}_{a \in \cO}$ be a collection of cone diagrams $\cR_a = \{p_{\alp}\colon K_{\alp}^{\triangleleft} \to \C_a\}$ in $\C_a$ for every $a \in \cO$. We then write $\Fun^{\cR}_{\cO}(\C,\D)^{\Day} \subseteq \Fun_{\cO}(\C,\D)^{\Day}$ for the full suboperad spanned by the $\cR_a$-local functors $f\colon \C_a \to \D_a$ for every $a$, and similarly $\Fun_{\cO}^{\cR-\mathrm{able}}(\C,\D)^{\Day} \subseteq \Fun_{\cO}(\C,\D)^{\Day}$ for the full suboperad spanned by the $\cR_a$-localisable functors for every $a \in \cO$. 

In this situation, we will say that $\cR$ is \defi{closed under $\cO$-tensoring} with objects in $\C$ if for every operation $\alp\colon (a_1,...,a_n) \to b$ in $\cO$ and objects $(c_2,\dots,c_n) \in \C_{a_2}\times \dots \times \C_{a_n}$, the functor $\alp(-,c_2,...,c_n)\colon \C_{a_1} \to \C_b$ sends diagrams in $\cR_{a_1}$ to diagrams in $\cR_b$. If $\C$ is in addition $\cO$-closed, we will say that $\cR$ is \defi{closed under $\cO$-cotensoring} with objects in $\C$ if for every operation $\alp\colon (a_1,...,a_n) \to b$ in $\cO$ and objects $c_i \in \C_{a_i}$ for $i=2,...n$, the right adjoint $R^{c_2,...,c_n}_{\alp}\colon \C_{b} \to \C_{a_1}$ sends diagrams in $\cR_{b}$ to diagrams in $\cR_{a_1}$.

\end{notation}

We now arrive at the first main result of this subsection.

\begin{proposition}
\label{proposition:multiplicative-day}%
Let $\C$ be an $\cO$-monoidal $\infty$-category
and let $\cR = \{\cR_a\}_a$ be a collection of cone diagrams $\cR_a = \{p_{\alp}\colon K_{\alpha}^{\triangleleft} \to \C_a\}_{\alp}$ for every $a \in \cO$. Let $\D$ be an almost closed $\cO$-monoidal $\infty$-category such that $\D_a$ admits $K_{\alp}$-indexed limits for every index $\alp$ and every $a \in \cO$ and suppose that $\cR$ is closed under $\cO$-tensoring with objects in $\C$.
Then the collection of left Bousfield localisations 
\[
\Fun_{\cO}^{\cR-\mathrm{able}}(\C,\D)_a = \Fun^{\cR-\mathrm{able}}(\C_a,\D_a) \locadjL{L_{a}} \Fun^{\cR}(\C_a,\D_a) = \Fun_{\cO}^{\cR}(\C,\D)_a
\]
for $a \in \cO$ is $\cO$-multiplicative in the sense of Definition~\reffour{definition:multiplicative-localisation} and so by Proposition~\reffour{proposition:multiplicative-localisation} assembles to form an $\cO$-multiplicative left Bousfield localisation of $\infty$-operads 
\[
\Fun_{\cO}^{\cR-\mathrm{able}}(\C,\D)^{\Day} \locadjL{L^{\otimes}} \Fun_{\cO}^{\cR}(\C,\D)^{\Day}.
\]
\end{proposition}


\begin{remark}
\label{remark:small-localisable}%
In the situation of Proposition~\reffour{proposition:multiplicative-day}, if $\D^{\otimes}$ is presentably $\cO$-monoidal and $\C$ is small 
then every functor $\C_a \to \D_a$ is $\cR_a$-localisable by the adjoint functor theorem and so, when $\cR$ is closed under $\cO$-tensoring with objects in $\C$, %
Proposition~\reffour{proposition:multiplicative-day} gives an $\cO$-monoidal left Bousfield localisation 
\[
\Fun_{\cO}(\C,\D)^{\Day} \locadjL{L^{\otimes}} \Fun_{\cO}^{\cR}(\C,\D)^\Day.
\]
of $\cO$-monoidal $\infty$-categories.
More generally, if $\C$ is accessible (and $\D^{\otimes}$ still presentably $\cO$-monoidal) and for every sufficiently large regular cardinal $\kappa$, every diagram in $\cR_a$ is a $\kappa$-filtered colimit of diagrams in $\cR_a$ involving only $\kappa$-compact objects, then the adjoint functor theorem implies that every small functor $f\colon \C_a \to \D_a$ admits an $\cR_a$-local approximation which is again small, and so the $\cO$-multiplicative left Bousfield localisation of Proposition~\reffour{proposition:multiplicative-day} restricts to an $\cO$-monoidal left Bousfield localisation 
\[
\Fun_{\cO}^{\sm}(\C,\D)^{\Day} \locadjL{L^{\otimes}_{\sm}} \Fun_{\cO}^{\cR,\sm}(\C,\D)^\Day
\]
on the level of small functors and small $\cR$-local functors, respectively, see Corollary~\reffour{corollary:persist-locally-small}. In this case, the collection of small functors also coincides with the collection of accessible ones.
\end{remark}



\begin{proof}[Proof of Proposition~\reffour{proposition:multiplicative-day}]
By Remark~\reffour{remark:full-suboperad}~\reffouritem{item:one-at-a-time} we need to show that for every operation $\alp\colon (a_1,...,a_n) \to b$ in $\cO$, any $\cR_{a_1}$-local equivalence $\sig\colon f_1 \Rightarrow f_1'$, any collection of functors $g_i\colon \C_{a_i} \to \D_{a_i}$ with $i=2,...,n$, and any $\cR_b$-local functor $h\colon \C_b \to \D_{b}$, the induced map
\[
\Nat(f_1' \boxtimes g_2 \boxtimes ... \boxtimes g_n,h|_{\C_{a_1} \times ... \times \C_{a_n}}) \to \Nat(f_1 \boxtimes g_2 \boxtimes ... \boxtimes g_n,h|_{\C_{a_1} \times ... \times \C_{a_n}})
\]
is an equivalence. We now observe that $\boxtimes$ preserves arbitrary colimits (not only small) in each entry separately by our assumption that $\D$ is almost closed. %
Writing each $g_i$ as a (possibly large) colimit of representables using Lemma~\reffour{lemma:canonical-colimit} we may reduce to the case where for each $i=2,\dots,n$ we have $g_i = (c_i)_!d_i$ for some $c_i \in \C$ and $d_i \in \D_{a_i}$. In addition, by the assumption that $\D^{\otimes} \to \cO^{\otimes}$ is almost closed
we may further assume that $d_i \in \D'_{a_i}$ (using the notation of Definition~\reffour{definition:almost-presentable}) for $i=2,...,n$, so that $\alp_{\D}(-,d_2,...,d_n)\colon \D_{a_1} \to \D_b$ admits a right adjoint $R^{d_2,...,d_n}_{\alp}\colon \D_b \to \D_{a_1}$. We then have
\begin{align*}
\Nat(f_1 \boxtimes (c_2)_!d_2 \boxtimes \dots \boxtimes (c_n)_!d_n, h|_{\C_{a_1}\times ... \times \C_{a_n}}) =& 
\Nat(\alp_{\D}(f_1,d_2,...,d_n), h(\alp_{\C}(-,c_2,...,c_n)))\\
=&\Nat(f_1,R^{d_2,...,d_n}_{\alp}(h(\alp_{\C}(-,c_2,...,c_n))))
\end{align*}
and similarly for $f_1'$.
Since $R^{d_2,...,d_n}_{\alp}$ preserves limits (being a right adjoint) it suffices to check that $h(\alp_{\C}(-,c_2,...,c_n))$
is $\cR_{a_1}$-local. This follows from the assumption that $\alp_{\C}(-,c_2,...,c_n)$ sends diagrams in $\cR_{a_1}$ to diagrams in $\cR_{b}$, and that $h$ is $\cR_b$-local.
\end{proof}




Combining Proposition~\reffour{proposition:multiplicative-day} with Remark~\reffour{remarks:symmetric-monoidal-objects-and-units}~\reffouritem{item:local-unit}
and Proposition~\reffour{proposition:units} %
we therefore conclude:


\begin{corollary}
\label{corollary:multiplicative-day}%
In the situation of Proposition~\reffour{proposition:multiplicative-day}, suppose in addition that $\cO$ admits a unit object $\emptyset \to \one_{\cO}$ (see Definition~\reffour{definition:monoidal-and-unit}), so that $\cB := \D_{\one_{\cO}}$ inherits a symmetric monoidal structure, and $\Fun_{\cO}(\C,\D)^{\Day}$ admits a unit $\one_{\Day} := (\one_\C)_!\one_{\cB}\colon \C_{\one_{\cO}} \to \D_{\one_{\cO}}$ over $\one_{\cO}$ (see Remark~\reffour{remark:units}). Assume that $\one_{\Day}$ is $\cR_{\one_{\cO}}$-localisable with $\cR_{\one_{\cO}}$-local approximation $\one_{\Day} \Rightarrow \one_\Day^{\cR}$. Then the composite $u\colon \emptyset \to \one_{\Day} \to \one_\Day^{\cR}$ exhibits $\one_\Day^{\cR}$ as a unit object in $\Fun_{\cO}^{\cR}(\C,\D)^{\Day}$ over $\one_{\cO}$
in the sense of Definition~\reffour{definition:monoidal-and-unit}. In particular, 
the functor 
\[
\one_\Day^{\cR}\colon \C_{\one_{\cO}} \to \D_{\one_{\cO}}
\]
carries an essentially unique lax symmetric monoidal structure with unit $u$ under which it is the initial $\cR$-local lax symmetric monoidal functor, and $\one_{\Day} \Rightarrow \one_\Day^{\cR}$ refines to a lax symmetric monoidal transformation in an essentially unique manner. In addition, for every $a \in \cO$, any $\cR_a$-local functor $\C_a \to \D_a$ carries an essentially unique $\one_\Day^\cR$-module structure. 
\end{corollary}


We now give a criterion for when the above multiplicative localisation is smashing.



\begin{proposition}
\label{proposition:smashing-day}%
Let $\cO^{\otimes}$ be an $\infty$-operad with a unit object $\emptyset \to \one_{\cO}$. 
Let $\C^{\otimes} \to \cO^{\otimes}$ be an $\cO$-closed $\cO$-monoidal $\infty$-category and let $\cR = \{\cR_a\}_a$ be a collection of cone diagrams $\cR_a = \{p_{\alp}\colon K_{\alpha}^{\triangleleft} \to \C_a\}_{\alp}$ for every $a \in \cO$. Let $\D^{\otimes} \to \cO^{\otimes}$ be an almost closed $\cO$-monoidal $\infty$-category such that $\D_a$ admits $K_{\alp}$-indexed limits for every index $\alp$ and every $a \in \cO$. In particular, $\C$ and $\D$ admits unit objects $\one_{\C} \in \C_{\one_{\cO}}$ and $\one_{\D} \in \D_{\one_{\cO}}$ over $\one_{\cO}$. Assume the following conditions are satisfied:
\begin{enumerate}
\item
$\cR$ is closed under $\cO$-tensoring and $\cO$-cotensoring with objects in $\C$ (see Notation~\reffour{notation:cone-diagrams}),
\item
\label{item:Rlocalclosedcolim}%
for every $a \in \cO$ the collection of $\cR_a$-local functors $\C_a \to \D_a$ is closed under all (potentially large) colimits that exist in $\Fun(\C_a,\D_a)$, 
and
\item
\label{item:OmonoidalpresKlimits}%
all the $\cO$-monoidal operations $\alp_{\D}\colon \D_{a_1} \times ... \times \D_{a_n} \to \D_b$ preserve, for each $i=1,...,n$, and each indexing $\infty$-category $K_{\alp}$ in $\cR_{a_i}$, $K_{\alp}$-indexed limits in the $i$'th argument.
\end{enumerate}
Then the following holds:
\begin{enumerate}
\item
\label{item:ideal}%
Let $a \in \cO$ and consider functors $F\colon \C_{\one_{\cO}} \to \D_{\one_{\cO}}$ and $G\colon \C_a \to \D_a$ at least one of whom is small, so that the Day convolution 
$F\otimes_{\Day} G\colon \C_a \to \D_a$ (covering the unique action of $\one_{\cO}$ on $a$) exists by Corollary~\reffour{corollary:persist-locally-small}\;\reffouritem{item:small-monoidal}. If at least one of $F,G$ is $\cR$-local then $F\otimes_{\Day} G$ is $\cR$-local.
\item
\label{item:equivalent}%
Let $\eta\colon\one_{\Day} = (\one_\C)_!\one_{\D} \Rightarrow F$ be a natural transformation of functors $\C_{\one_{\cO}} \to \D_{\one_{\cO}}$ such that $F$ is small and $\cR$-local. Then the following are equivalent:
\begin{enumerate}[label=\alph*),ref=\alph*)]
\item
\label{item:approximation}%
$\eta$ exhibits $F$ as an $\cR$-local approximation of $\one_{\Day}$.
\item
\label{item:localising}%
For every $a \in \cO$ and every $\cR$-local functor $G\colon \C_a \to \D_a$ the map $G \to F\otimes_{\Day} G$ is an equivalence.
\item
\label{item:F-smashing}%
The map $\eta$ exhibits $F$ as $\cO$-idempotent and for every $a \in \cO$ the collection of $\cR$-local functors $\C_a \to \D_a$ coincides with the image of $\F\otimes_{\Day}(-)\colon \Fun(\C_a,\D_a) \to \Fun(\C_a,\D_a)$.
\end{enumerate}
In particular, when these equivalent conditions hold all functors are $\cR$-localisable and
the resulting $\cO$-multiplicative left Bousfield localisation 
\[
\Fun_{\cO}(\C,\D)^{\Day} \locadjL{L^{\otimes}} \Fun_{\cO}^{\cR}(\C,\D)
\]
is $\cO$-smashing in the sense of Definition~\reffour{definition:smashing} with associated idempotent $F$.
\end{enumerate}
\end{proposition}




Under the same assumptions, the $\cO$-smashing localisation of Proposition~\reffour{proposition:smashing-day} restricts to an $\cO$-smashing localisation of $\cO$-monoidal $\infty$-categories 
\[
\Fun_{\cO}^{\sm}(\C,\D)^{\Day} \locadjL{L^{\otimes}_{\sm}} \Fun_{\cO}^{\cR,\sm}(\C,\D)
\]
on the level of small functors and small $\cR$-local functors, respectively, see Corollary~\reffour{corollary:persist-locally-small}.

We record the above special case of Proposition~\reffour{proposition:smashing-day}, which is frequently encountered:

\begin{corollary}
\label{corollary:smashing-day-small}%
Let $\C^{\otimes}$ be a small symmetric monoidal $\infty$-category, $\D^{\otimes}$ a presentably symmetric monoidal $\infty$-category and $\cR = \{p_{\alp}\colon K_{\alp}^{\triangleleft} \to \C\}_{\alp}$ a collection of cone diagrams in $\C$. Assume that the following conditions hold:
\begin{enumerate}
\item
the symmetric monoidal structure on $\C$ is closed,
\item
\label{item:tensor-and-cotensor}%
the collection $\cR$ is closed under both tensoring and cotensoring by objects in $\C$,
\item
\label{item:closed-colimits}%
the collection of $\cR$-local functors $\C \to \D$ is closed under all small colimits, 
and
\item
\label{item:closed-tensor}%
for $d \in \D$, the functor $d \otimes (-)\colon \D \to \D$ preserves $K_{\alp}$-indexed limits for every $K_{\alp}$ appearing in $\R$.
\end{enumerate}
Then the left Bousfield localisation $L\colon \Fun(\C,\D) \to \Fun^{\cR}(\C,\D)$ refines to a 
smashing left Bousfield localisation
\[
\Fun(\C,\D)^{\Day} \locadjL{L^{\otimes}} \Fun^{\cR}(\C,\D)^{\Day}
\]
of presentably symmetric monoidal $\infty$-categories.
\end{corollary}


\begin{example}
\label{example:excisive}%
Let $\C$ be a small stable $\infty$-category equipped with a closed symmetric monoidal structure and $\D$ a stable presentably symmetric monoidal $\infty$-category. Then tensoring and cotensoring with any object in $\C$ is automatically an exact operation and hence preserves for example, the collection of strongly cocartesian $n$-cubes. It then follows from Corollary~\reffour{corollary:smashing-day-small} that the full subcategory $\Exc^n(\C,\D)$ spanned by $n$-excisive functors is a smashing localisation of $\Fun(\C,\D)$, and similarly for the full subcategory $\Exc^n_*(\C,\D)$ of reduced $n$-excisive functors. For example, the $\infty$-category $\Funq(\C)$ of hermitian structures on $\C$ is a smashing localisation of $\Fun(\C\op,\Sp)$, a case of interest in the context of the present paper series.

For completeness, let us note that it is enough to assume above that the target $\D$ is closed and cocomplete, since the existence of an explicit $n$-excisive approximation is well-known in this case and does not require presentability, see~\cite{HA}*{Theorem 6.1.1.10}. We also point out that $\Exc^n(\C,\D)$ and $\Exc^n_*(\C,\D)$ are typically not smashing localisations of $\Fun(\C,\D)$ when $\C$ is not stable, see Example~\reffour{example:cant-remove} below.
\end{example}

\begin{remark}
\label{remark:rigid}%
Let $\C^{\otimes}$ be small symmetric monoidal $\infty$-category, $\D^{\otimes}$ a presentably symmetric monoidal $\infty$-category and $\cR$ a collection of cone diagrams in $\C$ which is closed under tensoring with objects in $\C$, so that $\Fun^{\cR}(\C,\D)$ is a multiplicative localisation of $\Fun(\C,\D)$ by Remark~\reffour{remark:small-localisable}. Now, if, for example, every object in $\C$ is dualisable then the monoidal structure on $\C$ is closed and $\cR$ is closed also under cotensoring, and similarly if $\D$ is generated under colimits by dualisable objects then Condition~\reffouritem{item:closed-colimits} implies Condition~\reffouritem{item:closed-tensor} from Corollary~\reffour{corollary:smashing-day-small}. In particular, when these dualisability assumptions hold for both $\C$ and $\D$, the conclusion of Corollary~\reffour{corollary:smashing-day-small} holds as soon as $\Fun^{\cR}(\C,\D)$ is closed under colimits. This is consistent with Remark~\reffour{remark:smashing-closed-colimits}: Indeed, in this case $\Fun^{\cR}(\C,\D)$ is itself generated under colimits by dualisable objects, and so 
its smashing localisations are exactly its multiplicative localisations which are closed under colimits. %
\end{remark}


\begin{example}
\label{example:cant-remove}%
In the situation of Corollary~\reffour{corollary:smashing-day-small}, the condition that $\cR$ is closed under cotensoring cannot be removed. For example, let $k$ be a field
and take
$\C = \Dperf_{\geq 0}(k)$ to be the symmetric monoidal $\infty$-category of connective perfect $k$-complexes, 
$\cR$ the collection of pushout squares in $\C$ together with the singleton diagram on the zero object, and $\D=\Spa$. Note that $\C$ is closed with internal mapping objects given by $\uline{\map}_{\geq 0}(P,Q) = \tau_{\geq 0}\hom_{\Dperf(k)}(P,Q)$.
Then all the conditions of Corollary~\reffour{corollary:smashing-day-small} are met except~\reffouritem{item:tensor-and-cotensor}, namely, $\cR$ is closed under tensoring but not under cotensoring. For example, cotensoring againt $k[1]$ %
doesn't preserve the cofibre sequence $k \to 0 \to k[1]$.
We claim that the resulting multiplicative localisation 
\[
\Fun(\Dperf(k)_{\geq 0},\Spa) \locadj \Fun^{\cR}(\Dperf(k)_{\geq 0},\Spa) = \Exc^1_*(\Dperf(k)_{\geq 0},\Spa)
\]
is not smashing. Indeed, for every $P \in \Dperf(k)_{\geq 0}$ the left adjoint sends $P_!\SS$ to $\map_{\Dperf(k)}(P,-)$, and in particular sends the unit $k \in \Dperf(k)_{\geq 0}$ to the underlying spectrum functor 
\[
U\colon \Dperf(k)_{\geq 0} \to \Sp .
\]
It hence suffices to show that the natural transformation
\[
U \otimes_{\Day} P_!\SS = U(\uline{\map}_{\geq 0}(P,-)) \Rightarrow \map_{\Dperf(k)}(P,-)
\]
is generally not an equivalence. %
Indeed, taking for example $P=k[1]$ and evaluating this transformation at $k$, the right hand side is $\map_{\Dperf(k)}(k[1],k) = k[-1]$ while the left hand side is $\tau_{\geq 0}k[-1] = 0$. %
\end{example}

\begin{proof}[Proof of Proposition~\reffour{proposition:smashing-day}]
We begin with proving~\reffouritem{item:ideal}. Assume first that $F$ is $\cR$-local.
Since the collection of $\cR_a$-local functors $\C_a \to \D_a$ is assumed to be closed under (potentially large) colimits and Day convolution preserves all colimits that exist in each entry since $\D$ is almost closed, it will suffice by Lemma~\reffour{lemma:canonical-colimit} to prove the claim for $G  = c_!d$ for some $c \in \C_a$ and $d \in \D_a$. As described in the proof of Proposition~\reffour{proposition:persist-locally-small},
the Day convolution $F \otimes_{\Day} c_!d$ is then given by $F(\uline{\Hom}_{\C_a}(c,-)) \otimes d$, which is $\cR_a$-local since $F$ is $\cR_{\one_{\cO}}$-local and we assumed that $\uline{\Hom}_{\C_a}(c,-)\colon \C_a \to \C_{\one_{\cO}}$ sends diagrams in $\cR_a$ to diagrams in $\cR_{\one_{\cO}}$ and that $d \otimes(-)\colon \D_{\one_{\cO}} \to \D_a$ preserves $K_{\alp}$-indexed limits for every $K_{\alp}$ in $\cR_{\one_{\cO}}$. The argument if $G$ is $\cR$-local is similar, where one now considers the case where $F=c_!d$ for some $c \in \C_{\one_{\cO}}$ and $d \in \D_{\one_{\cO}}$.

We now proceed to prove~\reffouritem{item:equivalent}. To see~\reffouritem{item:approximation}$\Rightarrow$\reffouritem{item:localising} note that if $\eta$ exhibits $F$ as an $\cR$-local approximation of $\one_{\Day}$ then by Proposition~\reffour{proposition:multiplicative-day} the map $G \Rightarrow F \otimes_{\Day} G$ is an $\cR_a$-local equivalence for every $G\colon \C_a \to \D_a$, and hence an equivalence for every $\cR$-local such $G$. Next, if \reffouritem{item:localising} holds then the induced map $F \to F\otimes_{\Day} F$ is an equivalence since $F$ is $\cR$-local, and so $\eta$ exhibits $F$ is as $\cO$-idempotent. In addition, \reffouritem{item:localising} implies that the the image of $\F\otimes_{\Day}(-)\colon \Fun(\C_a,\D_a) \to \Fun(\C_a,\D_a)$ contains all $\cR$-local functors and the inverse inclusion follows from~\reffouritem{item:ideal}. We thus have \reffouritem{item:localising}$\Rightarrow$\reffouritem{item:F-smashing}. Finally the implication \reffouritem{item:F-smashing}$\Rightarrow$\reffouritem{item:approximation} follows from Remark~\reffour{remark:action-of-idempotent}.

The last part of~\reffouritem{item:equivalent} now follows from Proposition~\reffour{proposition:equivalent-1}.
\end{proof}




Let us finish with a few comments on the conditions on the target $\infty$-category in Proposition~\reffour{proposition:smashing-day}.

\begin{examples}
\label{someexampleofRlocalfunctors}%
Let $\cO^{\otimes},\C^{\otimes}$ and $\cR$ be as in Proposition~\reffour{proposition:smashing-day}.
\begin{enumerate}
\item If each diagram in each $\cR_a$ has source $(\Delta^0)^\triangleleft = \Delta^1$, then each $\cR_a$ is simply a collection of arrows in $\C_a$ and a functor on $\C_a$ %
is $\cR_a$-local if and only if it inverts all these maps. Then any $\cO$-monoidal $\D^{\otimes}$ satisfies assumptions \reffouritem{item:Rlocalclosedcolim} and \reffouritem{item:OmonoidalpresKlimits} of Proposition \reffour{proposition:smashing-day}.
\item If each diagram in each $\cR_a$ has source $\emptyset^\triangleleft = \Delta^0$, then $\cR_a$ is simply a collection of objects in $\C_a$ and a functor on $\C_a$ is $\cR_a$-local if and only it takes all these objects to the terminal object. %
Then any pointedly $\cO$-monoidal $\D$ satisfies assumptions \reffouritem{item:Rlocalclosedcolim} and \reffouritem{item:OmonoidalpresKlimits}. %
\item If all sources of morphisms in each $\cR_a$ are cones on finite sets, then any semi-additivly $\cO$-monoidal $\D$ satisfies assumptions \reffouritem{item:Rlocalclosedcolim} and \reffouritem{item:OmonoidalpresKlimits}. %
Similarly, if the sources are all finite $\infty$-categories then the same holds for any stably $\cO$-monoidal $\D^{\otimes}$.
\item As a slightly more subtle example, observe that if $\cR_a$ consists of square diagrams whose vertical maps admit sections, then any pre-stable $\D$ can serve as the target, since squares in $\D$ which admit vertical sections are pullbacks if and only if they are pushouts. 

A relevant example for our paper series is $\C^{\otimes} = (\Catx)^{\otimes} \to \Com^{\otimes}$, $\cR$ the collection of split Verdier squares and $\D = \Grp_{\Einf}$, in which case $\one_\Day^{\cR} = \Kspace$.
\end{enumerate}
\end{examples}
